\documentclass[8pt]{beamer}
\usetheme{Madrid}
\usecolortheme{default}
\usepackage{amsmath, amssymb, graphicx, array, multirow, booktabs, xcolor, tikz}
\usetikzlibrary{shapes,arrows,positioning,fit,backgrounds}

% Compact font settings
\setbeamerfont{title}{size=\Large}
\setbeamerfont{subtitle}{size=\small}
\setbeamerfont{author}{size=\scriptsize}
\setbeamerfont{date}{size=\scriptsize}
\setbeamerfont{frametitle}{size=\large}
\setbeamerfont{framesubtitle}{size=\normalsize}
\setbeamerfont{enumerate item}{size=\footnotesize}
\setbeamerfont{itemize item}{size=\footnotesize}
\setbeamerfont{block title}{size=\footnotesize}

% Compact spacing
\setlength{\itemsep}{0.08ex}
\setlength{\parskip}{0.08ex}
\setlength{\leftmargini}{0.3cm}
\setbeamersize{text margin left=3mm, text margin right=3mm}
\addtobeamertemplate{frame begin}{\vspace{-0.4cm}}{}

% Colors
\definecolor{myblue}{RGB}{0,0,255}
\definecolor{myred}{RGB}{255,0,0}
\definecolor{mygreen}{RGB}{0,128,0}
\definecolor{myorange}{RGB}{255,165,0}
\definecolor{mypurple}{RGB}{128,0,128}
\definecolor{mygray}{RGB}{128,128,128}
\definecolor{mygold}{RGB}{255,215,0}

\title{Responsible \& Ethical AI}
\subtitle{100 Questions: Ethical Frameworks, Principles, Bias, Governance, Regulation}
\author{GARP RAI Exam Preparation}
\date{}

\begin{document}
	
	\begin{frame}
		\titlepage
	\end{frame}
	
	% ============================================================
	% SECTION 1: INTRODUCTION TO RESPONSIBLE AI - QUESTIONS 1-5
	% ============================================================
	
	\begin{frame}{Q1: Responsible AI Definition - Tutorial}
		\fontsize{6.5}{7.5}\selectfont
		\textbf{QUESTION: What is Responsible AI?}
		
		\vspace{0.1cm}
		\textbf{OPTIONS:}
		\begin{enumerate}
			\item A marketing term with no practical meaning
			\item The ethical and responsible development, deployment, and use of AI systems to benefit society while mitigating risks
			\item A specific type of machine learning algorithm
			\item A regulatory framework only applicable in Europe
		\end{enumerate}
		
		\vspace{0.1cm}
		\textcolor{myblue}{\textbf{ANSWER: B}}
		
		\vspace{0.1cm}
		\textbf{DETAILED EXPLANATION:}
		\begin{itemize}
			\item \textbf{Why B is correct:} Responsible AI aims to ensure AI benefits society holistically while mitigating potential risks. It involves aligning AI systems with ethical standards, industry standards, regulation, legislation, and global principles.
			\item \textbf{Why A is wrong:} Responsible AI has substantive practical meaning and applications
			\item \textbf{Why C is wrong:} It's a governance and ethics framework, not a specific algorithm
			\item \textbf{Why D is wrong:} Responsible AI is a global concept, not limited to Europe
		\end{itemize}
		
		\textcolor{myred}{\textbf{EXAM TRAP:}} Responsible AI = \textcolor{myblue}{ethical development and deployment of AI systems}!
	\end{frame}
	
	\begin{frame}{Q2: Ethical Frameworks Purpose - Tutorial}
		\fontsize{6.5}{7.5}\selectfont
		\textbf{QUESTION: What is the purpose of ethical frameworks in AI?}
		
		\vspace{0.1cm}
		\textbf{OPTIONS:}
		\begin{enumerate}
			\item To make AI development slower
			\item To help decide how to design and deploy AI to minimize risks, harms, and wrongs
			\item To replace all technical considerations
			\item To eliminate all risks completely
		\end{enumerate}
		
		\vspace{0.1cm}
		\textcolor{myblue}{\textbf{ANSWER: B}}
		
		\vspace{0.1cm}
		\textbf{DETAILED EXPLANATION:}
		\begin{itemize}
			\item \textbf{Why B is correct:} Ethical frameworks help define principles to abide by, identify unacceptable biases, assess transparency needs, identify stakeholders, promote accountability, and make justifiable decisions.
			\item \textbf{Why A is wrong:} Ethics aims to guide, not slow, development
			\item \textbf{Why C is wrong:} Technical considerations remain essential
			\item \textbf{Why D is wrong:} Complete risk elimination is impossible
		\end{itemize}
		
		\textcolor{myred}{\textbf{EXAM TRAP:}} Ethical frameworks \textcolor{myblue}{guide decision-making} to minimize risks and harms!
	\end{frame}
	
	\begin{frame}{Q3: Pluralistic Societies and Ethics - Tutorial}
		\fontsize{6.5}{7.5}\selectfont
		\textbf{QUESTION: Why are different ethical frameworks important in pluralistic societies?}
		
		\vspace{0.1cm}
		\textbf{OPTIONS:}
		\begin{enumerate}
			\item Because everyone agrees on ethics
			\item Because they help us make justifiable decisions when people hold different values
			\item Because only one framework is ever needed
			\item Because ethics is not important in diverse societies
		\end{enumerate}
		
		\vspace{0.1cm}
		\textcolor{myblue}{\textbf{ANSWER: B}}
		
		\vspace{0.1cm}
		\textbf{DETAILED EXPLANATION:}
		\begin{itemize}
			\item \textbf{Why B is correct:} Different ethical frameworks help us make justifiable decisions in pluralistic societies where people hold diverse values and beliefs.
			\item \textbf{Why A is wrong:} People have different ethical views
			\item \textbf{Why C is wrong:} Multiple frameworks provide complementary insights
			\item \textbf{Why D is wrong:} Ethics is particularly important in diverse societies
		\end{itemize}
		
		\textcolor{myred}{\textbf{EXAM TRAP:}} Different frameworks help make \textcolor{myblue}{justifiable decisions} in diverse societies!
	\end{frame}
	
	\begin{frame}{Q4: Practical Ethics Definition - Tutorial}
		\fontsize{6.5}{7.5}\selectfont
		\textbf{QUESTION: What is practical ethics?}
		
		\vspace{0.1cm}
		\textbf{OPTIONS:}
		\begin{enumerate}
			\item Theoretical philosophy disconnected from reality
			\item The application of ethical theory and development of good practices to solve real-world moral dilemmas
			\item A legal framework for AI regulation
			\item A type of machine learning algorithm
		\end{enumerate}
		
		\vspace{0.1cm}
		\textcolor{myblue}{\textbf{ANSWER: B}}
		
		\vspace{0.1cm}
		\textbf{DETAILED EXPLANATION:}
		\begin{itemize}
			\item \textbf{Why B is correct:} Practical ethics applies ethical theory to solve real-world moral dilemmas, providing concrete guidance for moral decision making.
			\item \textbf{Why A is wrong:} It is applied, not purely theoretical
			\item \textbf{Why C is wrong:} It's broader than law, informing legal frameworks
			\item \textbf{Why D is wrong:} It's about ethics, not algorithms
		\end{itemize}
		
		\textcolor{myred}{\textbf{EXAM TRAP:}} Practical ethics = \textcolor{myblue}{applying theory to real-world dilemmas}!
	\end{frame}
	
	\begin{frame}{Q5: Practical Ethics Characteristics - Tutorial}
		\fontsize{6.5}{7.5}\selectfont
		\textbf{QUESTION: Which is NOT a characteristic of practical ethics?}
		
		\vspace{0.1cm}
		\textbf{OPTIONS:}
		\begin{enumerate}
			\item Applied focus on actual ethical dilemmas
			\item Multidisciplinary approach drawing on various fields
			\item Context-sensitive emphasis on situational nuances
			\item Ignoring real-world consequences
		\end{enumerate}
		
		\vspace{0.1cm}
		\textcolor{myblue}{\textbf{ANSWER: D}}
		
		\vspace{0.1cm}
		\textbf{DETAILED EXPLANATION:}
		\begin{itemize}
			\item \textbf{Why D is correct:} Practical ethics explicitly focuses on real-world consequences and providing concrete guidance, not ignoring them.
			\item \textbf{Why A is wrong:} Applied focus IS a characteristic
			\item \textbf{Why B is wrong:} Multidisciplinary approach IS a characteristic
			\item \textbf{Why C is wrong:} Context sensitivity IS a characteristic
		\end{itemize}
		
		\textcolor{myred}{\textbf{EXAM TRAP:}} Practical ethics is \textcolor{myblue}{applied, multidisciplinary, and context-sensitive} - never ignores consequences!
	\end{frame}
	
	% ============================================================
	% SECTION 2: BENEFITS OF ETHICAL FRAMEWORKS - QUESTIONS 6-15
	% ============================================================
	
	\begin{frame}{Q6: Building Trust and Reputation - Tutorial}
		\fontsize{6.5}{7.5}\selectfont
		\textbf{QUESTION: How does a practical ethics framework help build trust and reputation?}
		
		\vspace{0.1cm}
		\textbf{OPTIONS:}
		\begin{enumerate}
			\item By making companies look good without real change
			\item Customers, employees, and the public are more likely to trust companies known for adhering to ethical standards
			\item By hiding unethical practices
			\item By avoiding all public scrutiny
		\end{enumerate}
		
		\vspace{0.1cm}
		\textcolor{myblue}{\textbf{ANSWER: B}}
		
		\vspace{0.1cm}
		\textbf{DETAILED EXPLANATION:}
		\begin{itemize}
			\item \textbf{Why B is correct:} Companies known for ethical standards build trust with customers, employees, partners, and the public, improving brand image and reputation.
			\item \textbf{Why A is wrong:} Ethics requires genuine commitment, not just appearance
			\item \textbf{Why C is wrong:} Hiding practices undermines trust
			\item \textbf{Why D is wrong:} Public scrutiny is part of accountability
		\end{itemize}
		
		\textcolor{myred}{\textbf{EXAM TRAP:}} Ethics frameworks build \textcolor{myblue}{trust and reputation} through genuine adherence to standards!
	\end{frame}
	
	\begin{frame}{Q7: Avoiding Scandals - Tutorial}
		\fontsize{6.5}{7.5}\selectfont
		\textbf{QUESTION: How can a practical ethics framework help avoid scandals?}
		
		\vspace{0.1cm}
		\textbf{OPTIONS:}
		\begin{enumerate}
			\item By ignoring potential risks
			\item Unethical behavior like fraud can damage reputation; ethics helps prevent major scandals
			\item By making scandals less visible
			\item By blaming employees for failures
		\end{enumerate}
		
		\vspace{0.1cm}
		\textcolor{myblue}{\textbf{ANSWER: B}}
		
		\vspace{0.1cm}
		\textbf{DETAILED EXPLANATION:}
		\begin{itemize}
			\item \textbf{Why B is correct:} Unethical behavior (fraud, bias, privacy violations) damages reputation and bottom line. Practical ethics helps prevent major scandals by embedding ethical practices.
			\item \textbf{Why A is wrong:} Ignoring risks increases scandal likelihood
			\item \textbf{Why C is wrong:} Scandals become visible regardless
			\item \textbf{Why D is wrong:} Ethics frameworks promote accountability, not scapegoating
		\end{itemize}
		
		\textcolor{myred}{\textbf{EXAM TRAP:}} Ethics helps \textcolor{myblue}{prevent scandals} by embedding ethical practices!
	\end{frame}
	
	\begin{frame}{Q8: Attracting and Retaining Talent - Tutorial}
		\fontsize{6.5}{7.5}\selectfont
		\textbf{QUESTION: How does a practical ethics framework help with talent?}
		
		\vspace{0.1cm}
		\textbf{OPTIONS:}
		\begin{enumerate}
			\item By paying higher salaries
			\item Workers increasingly care about responsible and ethical practices, so ethics frameworks help recruit and retain talent
			\item By reducing workforce requirements
			\item By outsourcing ethical decisions
		\end{enumerate}
		
		\vspace{0.1cm}
		\textcolor{myblue}{\textbf{ANSWER: B}}
		
		\vspace{0.1cm}
		\textbf{DETAILED EXPLANATION:}
		\begin{itemize}
			\item \textbf{Why B is correct:} Employees entering the job market are increasingly concerned about their companies' ethical practices. Having a practical ethics framework helps attract and retain talent.
			\item \textbf{Why A is wrong:} Ethics is about values, not just compensation
			\item \textbf{Why C is wrong:} Ethics frameworks don't reduce workforce needs
			\item \textbf{Why D is wrong:} Ethics decisions cannot be fully outsourced
		\end{itemize}
		
		\textcolor{myred}{\textbf{EXAM TRAP:}} Ethics frameworks help \textcolor{myblue}{attract and retain talent} concerned about responsible practices!
	\end{frame}
	
	\begin{frame}{Q9: Strengthening Culture - Tutorial}
		\fontsize{6.5}{7.5}\selectfont
		\textbf{QUESTION: How does a practical ethics framework strengthen workplace culture?}
		
		\vspace{0.1cm}
		\textbf{OPTIONS:}
		\begin{enumerate}
			\item By eliminating all disagreements
			\item Ethics programs can nurture teamwork, accountability, and integrity, boosting morale and productivity
			\item By enforcing strict rules without flexibility
			\item By ignoring employee concerns
		\end{enumerate}
		
		\vspace{0.1cm}
		\textcolor{myblue}{\textbf{ANSWER: B}}
		
		\vspace{0.1cm}
		\textbf{DETAILED EXPLANATION:}
		\begin{itemize}
			\item \textbf{Why B is correct:} Ethics programs nurture teamwork, accountability, and integrity in workplace culture, which in turn boosts morale and productivity.
			\item \textbf{Why A is wrong:} Disagreements are natural and can be healthy
			\item \textbf{Why C is wrong:} Flexibility is important in ethics
			\item \textbf{Why D is wrong:} Employee concerns should be addressed
		\end{itemize}
		
		\textcolor{myred}{\textbf{EXAM TRAP:}} Ethics frameworks strengthen \textcolor{myblue}{teamwork, accountability, and integrity} in workplace culture!
	\end{frame}
	
	\begin{frame}{Q10: Supporting Risk Management - Tutorial}
		\fontsize{6.5}{7.5}\selectfont
		\textbf{QUESTION: How does a practical ethics framework support risk management?}
		
		\vspace{0.1cm}
		\textbf{OPTIONS:}
		\begin{enumerate}
			\item By ignoring all risks
			\item Ethical frameworks can help companies reduce risk of safety failures, data breaches, and regulatory non-compliance
			\item By transferring all risk to customers
			\item By avoiding regulation altogether
		\end{enumerate}
		
		\vspace{0.1cm}
		\textcolor{myblue}{\textbf{ANSWER: B}}
		
		\vspace{0.1cm}
		\textbf{DETAILED EXPLANATION:}
		\begin{itemize}
			\item \textbf{Why B is correct:} Having an ethical framework undergirding policies and procedures can help companies reduce risk of safety failures, data breaches, and regulatory non-compliance.
			\item \textbf{Why A is wrong:} Ignoring risks increases exposure
			\item \textbf{Why C is wrong:} Risk transfer is limited and often unethical
			\item \textbf{Why D is wrong:} Regulation is a reality that must be addressed
		\end{itemize}
		
		\textcolor{myred}{\textbf{EXAM TRAP:}} Ethics frameworks \textcolor{myblue}{reduce risks} of failures, breaches, and non-compliance!
	\end{frame}
	
	\begin{frame}{Q11: Guiding Innovation - Tutorial}
		\fontsize{6.5}{7.5}\selectfont
		\textbf{QUESTION: How does an ethical framework guide innovation?}
		
		\vspace{0.1cm}
		\textbf{OPTIONS:}
		\begin{enumerate}
			\item By stopping all innovation
			\item By helping guide innovation in a responsible direction that considers impacts on people and society
			\item By making innovation more expensive
			\item By ignoring societal impacts
		\end{enumerate}
		
		\vspace{0.1cm}
		\textcolor{myblue}{\textbf{ANSWER: B}}
		
		\vspace{0.1cm}
		\textbf{DETAILED EXPLANATION:}
		\begin{itemize}
			\item \textbf{Why B is correct:} An ethical framework helps guide innovation in a responsible direction that considers impacts on people and society, ensuring progress benefits everyone.
			\item \textbf{Why A is wrong:} Ethics enables, not stops, innovation
			\item \textbf{Why C is wrong:} While there may be costs, the benefits outweigh them
			\item \textbf{Why D is wrong:} Societal impacts are central to ethics
		\end{itemize}
		
		\textcolor{myred}{\textbf{EXAM TRAP:}} Ethics guides innovation in \textcolor{myblue}{responsible directions} considering societal impacts!
	\end{frame}
	
	\begin{frame}{Q12: Promoting Long-Term Thinking - Tutorial}
		\fontsize{6.5}{7.5}\selectfont
		\textbf{QUESTION: How does ethics promote long-term thinking?}
		
		\vspace{0.1cm}
		\textbf{OPTIONS:}
		\begin{enumerate}
			\item By focusing only on short-term profits
			\item Managing in an environment with ethical considerations helps promote a long-term view
			\item By avoiding any future planning
			\item By prioritizing quarterly results
		\end{enumerate}
		
		\vspace{0.1cm}
		\textcolor{myblue}{\textbf{ANSWER: B}}
		
		\vspace{0.1cm}
		\textbf{DETAILED EXPLANATION:}
		\begin{itemize}
			\item \textbf{Why B is correct:} Managing in an environment with ethical considerations helps promote a long-term view, considering sustainability and future impacts.
			\item \textbf{Why A is wrong:} Ethics balances short and long-term considerations
			\item \textbf{Why C is wrong:} Ethics requires planning for the future
			\item \textbf{Why D is wrong:} Ethics transcends quarterly results
		\end{itemize}
		
		\textcolor{myred}{\textbf{EXAM TRAP:}} Ethics promotes \textcolor{myblue}{long-term thinking} and sustainability!
	\end{frame}
	
	\begin{frame}{Q13: Business Case for Ethics - Tutorial}
		\fontsize{6.5}{7.5}\selectfont
		\textbf{QUESTION: What is the business case for practical ethics?}
		
		\vspace{0.1cm}
		\textbf{OPTIONS:}
		\begin{enumerate}
			\item Ethics always reduces profits
			\item Proponents argue that practical ethics makes good business sense by providing multiple benefits
			\item Ethics is only about avoiding fines
			\item Ethics has no business value
		\end{enumerate}
		
		\vspace{0.1cm}
		\textcolor{myblue}{\textbf{ANSWER: B}}
		
		\vspace{0.1cm}
		\textbf{DETAILED EXPLANATION:}
		\begin{itemize}
			\item \textbf{Why B is correct:} Practical ethics provides multiple business benefits: trust building, scandal avoidance, talent attraction, culture strengthening, risk management, innovation guidance, and long-term thinking.
			\item \textbf{Why A is wrong:} Ethics can enhance long-term profitability
			\item \textbf{Why C is wrong:} Ethics is broader than just avoiding fines
			\item \textbf{Why D is wrong:} Ethics has significant business value
		\end{itemize}
		
		\textcolor{myred}{\textbf{EXAM TRAP:}} Ethics makes \textcolor{myblue}{good business sense} with multiple benefits!
	\end{frame}
	
	\begin{frame}{Q14: Ethics and Regulation Relationship - Tutorial}
		\fontsize{6.5}{7.5}\selectfont
		\textbf{QUESTION: How do ethics and law relate?}
		
		\vspace{0.1cm}
		\textbf{OPTIONS:}
		\begin{enumerate}
			\item Ethics and law are the same thing
			\item Ethics is a complement to law, helping ground, inform, and shape law
			\item Ethics is irrelevant to law
			\item Law always reflects ethics perfectly
		\end{enumerate}
		
		\vspace{0.1cm}
		\textcolor{myblue}{\textbf{ANSWER: B}}
		
		\vspace{0.1cm}
		\textbf{DETAILED EXPLANATION:}
		\begin{itemize}
			\item \textbf{Why B is correct:} Ethics complements law and is necessary to ground, inform, and shape law. Ethics helps distinguish between just and unjust laws.
			\item \textbf{Why A is wrong:} They are distinct but related
			\item \textbf{Why C is wrong:} Ethics informs law
			\item \textbf{Why D is wrong:} Law often lags behind ethical understanding
		\end{itemize}
		
		\textcolor{myred}{\textbf{EXAM TRAP:}} Ethics \textcolor{myblue}{complements and informs} law!
	\end{frame}
	
	\begin{frame}{Q15: Ethics vs Law Scope - Tutorial}
		\fontsize{6.5}{7.5}\selectfont
		\textbf{QUESTION: How does ethics differ from law in scope?}
		
		\vspace{0.1cm}
		\textbf{OPTIONS:}
		\begin{enumerate}
			\item Ethics is narrower than law
			\item Laws establish minimal requirements; ethics goes beyond, identifying moral issues and making recommendations
			\item Ethics only applies to individuals, not organizations
			\item Law covers all ethical concerns
		\end{enumerate}
		
		\vspace{0.1cm}
		\textcolor{myblue}{\textbf{ANSWER: B}}
		
		\vspace{0.1cm}
		\textbf{DETAILED EXPLANATION:}
		\begin{itemize}
			\item \textbf{Why B is correct:} Laws establish minimal behavioral requirements. Ethics goes beyond, identifying moral issues, reflecting on the kind of society we want, and making recommendations for a good life.
			\item \textbf{Why A is wrong:} Ethics is broader than law
			\item \textbf{Why C is wrong:} Ethics applies to organizations too
			\item \textbf{Why D is wrong:} Law doesn't cover all ethical concerns
		\end{itemize}
		
		\textcolor{myred}{\textbf{EXAM TRAP:}} Ethics is \textcolor{myblue}{more ambitious than law} - establishes ideals, not just minimums!
	\end{frame}
	
	% ============================================================
	% SECTION 3: ETHICAL FRAMEWORKS - QUESTIONS 16-35
	% ============================================================
	
	\begin{frame}{Q16: Consequentialism Definition - Tutorial}
		\fontsize{6.5}{7.5}\selectfont
		\textbf{QUESTION: What is consequentialism?}
		
		\vspace{0.1cm}
		\textbf{OPTIONS:}
		\begin{enumerate}
			\item An ethical theory that judges morality based on adherence to rules
			\item An ethical theory that judges morality based on the consequences of actions
			\item An ethical theory that focuses on virtuous character traits
			\item An ethical theory that ignores all outcomes
		\end{enumerate}
		
		\vspace{0.1cm}
		\textcolor{myblue}{\textbf{ANSWER: B}}
		
		\vspace{0.1cm}
		\textbf{DETAILED EXPLANATION:}
		\begin{itemize}
			\item \textbf{Why B is correct:} Consequentialism judges the morality of an action based on its consequences. The most common form is utilitarianism, which aims to maximize overall utility.
			\item \textbf{Why A is wrong:} This is deontology
			\item \textbf{Why C is wrong:} This is virtue ethics
			\item \textbf{Why D is wrong:} Consequentialism focuses heavily on outcomes
		\end{itemize}
		
		\textcolor{myred}{\textbf{EXAM TRAP:}} Consequentialism = \textcolor{myblue}{judging actions by their consequences}!
	\end{frame}
	
	\begin{frame}{Q17: Utilitarianism Definition - Tutorial}
		\fontsize{6.5}{7.5}\selectfont
		\textbf{QUESTION: What is utilitarianism?}
		
		\vspace{0.1cm}
		\textbf{OPTIONS:}
		\begin{enumerate}
			\item The most common form of consequentialism that aims to maximize overall utility
			\item A deontological theory focused on duties
			\item A virtue ethics approach to character
			\item A theory that ignores consequences
		\end{enumerate}
		
		\vspace{0.1cm}
		\textcolor{myblue}{\textbf{ANSWER: A}}
		
		\vspace{0.1cm}
		\textbf{DETAILED EXPLANATION:}
		\begin{itemize}
			\item \textbf{Why A is correct:} Utilitarianism is the most common form of consequentialism. It aims to maximize overall utility, often defined in terms of pleasure, happiness, or satisfaction of desires.
			\item \textbf{Why B is wrong:} Utilitarianism is consequentialist, not deontological
			\item \textbf{Why C is wrong:} Utilitarianism focuses on consequences, not character
			\item \textbf{Why D is wrong:} Consequences are central to utilitarianism
		\end{itemize}
		
		\textcolor{myred}{\textbf{EXAM TRAP:}} Utilitarianism = \textcolor{myblue}{maximize overall utility/happiness}!
	\end{frame}
	
	\begin{frame}{Q18: Utilitarian Thinkers - Tutorial}
		\fontsize{6.5}{7.5}\selectfont
		\textbf{QUESTION: Who were two famous proponents of utilitarianism?}
		
		\vspace{0.1cm}
		\textbf{OPTIONS:}
		\begin{enumerate}
			\item Plato and Aristotle
			\item Jeremy Bentham and John Stuart Mill
			\item Immanuel Kant and Aristotle
			\item Confucius and Buddha
		\end{enumerate}
		
		\vspace{0.1cm}
		\textcolor{myblue}{\textbf{ANSWER: B}}
		
		\vspace{0.1cm}
		\textbf{DETAILED EXPLANATION:}
		\begin{itemize}
			\item \textbf{Why B is correct:} Jeremy Bentham and John Stuart Mill are the two most famous proponents of utilitarianism, developing the theory of maximizing utility.
			\item \textbf{Why A is wrong:} Plato and Aristotle are associated with virtue ethics
			\item \textbf{Why C is wrong:} Kant is associated with deontology
			\item \textbf{Why D is wrong:} These are religious/spiritual figures, not utilitarians
		\end{itemize}
		
		\textcolor{myred}{\textbf{EXAM TRAP:}} Key utilitarians = \textcolor{myblue}{Bentham and Mill}!
	\end{frame}
	
	\begin{frame}{Q19: Consequentialism Advantage - Tutorial}
		\fontsize{6.5}{7.5}\selectfont
		\textbf{QUESTION: What is a key advantage of consequentialism?}
		
		\vspace{0.1cm}
		\textbf{OPTIONS:}
		\begin{enumerate}
			\item It provides a single, quantifiable metric for determining moral value
			\item It never leads to difficult decisions
			\item It always produces the right answer
			\item It ignores all values
		\end{enumerate}
		
		\vspace{0.1cm}
		\textcolor{myblue}{\textbf{ANSWER: A}}
		
		\vspace{0.1cm}
		\textbf{DETAILED EXPLANATION:}
		\begin{itemize}
			\item \textbf{Why A is correct:} Consequentialism provides a single, quantifiable metric for determining moral value, making decision-making more straightforward.
			\item \textbf{Why B is wrong:} Consequentialism still involves difficult decisions
			\item \textbf{Why C is wrong:} No ethical theory always produces the right answer
			\item \textbf{Why D is wrong:} It's based on values of utility/happiness
		\end{itemize}
		
		\textcolor{myred}{\textbf{EXAM TRAP:}} Consequentialism advantage = \textcolor{myblue}{single quantifiable metric}!
	\end{frame}
	
	\begin{frame}{Q20: Consequentialism Criticism - Tutorial}
		\fontsize{6.5}{7.5}\selectfont
		\textbf{QUESTION: What is a criticism of consequentialism?}
		
		\vspace{0.1cm}
		\textbf{OPTIONS:}
		\begin{enumerate}
			\item It is too focused on rules
			\item It can justify violating individual rights for the greater good
			\item It ignores consequences
			\item It is too focused on character
		\end{enumerate}
		
		\vspace{0.1cm}
		\textcolor{myblue}{\textbf{ANSWER: B}}
		
		\vspace{0.1cm}
		\textbf{DETAILED EXPLANATION:}
		\begin{itemize}
			\item \textbf{Why B is correct:} Critics argue that maximizing utility can lead to actions many consider immoral, like severely violating individual rights for the greater good (e.g., killing one person to save five).
			\item \textbf{Why A is wrong:} Rules are more central to deontology
			\item \textbf{Why C is wrong:} Consequentialism focuses heavily on consequences
			\item \textbf{Why D is wrong:} Character is central to virtue ethics
		\end{itemize}
		
		\textcolor{myred}{\textbf{EXAM TRAP:}} Consequentialism can justify \textcolor{myblue}{violating rights for the greater good}!
	\end{frame}
	
	\begin{frame}{Q21: Rule Consequentialism - Tutorial}
		\fontsize{6.5}{7.5}\selectfont
		\textbf{QUESTION: What is rule consequentialism?}
		
		\vspace{0.1cm}
		\textbf{OPTIONS:}
		\begin{enumerate}
			\item Judging acts by whether they adhere to utility-maximizing rules
			\item Always breaking rules for better outcomes
			\item Following all rules regardless of consequences
			\item Ignoring rules completely
		\end{enumerate}
		
		\vspace{0.1cm}
		\textcolor{myblue}{\textbf{ANSWER: A}}
		
		\vspace{0.1cm}
		\textbf{DETAILED EXPLANATION:}
		\begin{itemize}
			\item \textbf{Why A is correct:} Rule consequentialists judge acts by whether they adhere to utility-maximizing rules, not only by their case-specific outcomes. This addresses issues with act utilitarianism.
			\item \textbf{Why B is wrong:} This is act utilitarianism, not rule consequentialism
			\item \textbf{Why C is wrong:} This is deontology, not consequentialism
			\item \textbf{Why D is wrong:} This ignores ethical considerations
		\end{itemize}
		
		\textcolor{myred}{\textbf{EXAM TRAP:}} Rule consequentialism = \textcolor{myblue}{follow utility-maximizing rules}!
	\end{frame}
	
	\begin{frame}{Q22: Deontology Definition - Tutorial}
		\fontsize{6.5}{7.5}\selectfont
		\textbf{QUESTION: What is deontology?}
		
		\vspace{0.1cm}
		\textbf{OPTIONS:}
		\begin{enumerate}
			\item An ethical theory that judges morality based on consequences
			\item An ethical framework that judges morality based on adherence to ethical duties and rules
			\item An ethical theory that focuses on virtuous character traits
			\item An ethical theory that ignores all moral considerations
		\end{enumerate}
		
		\vspace{0.1cm}
		\textcolor{myblue}{\textbf{ANSWER: B}}
		
		\vspace{0.1cm}
		\textbf{DETAILED EXPLANATION:}
		\begin{itemize}
			\item \textbf{Why B is correct:} Deontology judges morality based on adherence to ethical duties and rules rather than focusing on consequences. It emphasizes moral obligations regardless of outcomes.
			\item \textbf{Why A is wrong:} This is consequentialism
			\item \textbf{Why C is wrong:} This is virtue ethics
			\item \textbf{Why D is wrong:} Deontology is based on moral considerations
		\end{itemize}
		
		\textcolor{myred}{\textbf{EXAM TRAP:}} Deontology = \textcolor{myblue}{judging actions by duties and rules}!
	\end{frame}
	
	\begin{frame}{Q23: Kant and Deontology - Tutorial}
		\fontsize{6.5}{7.5}\selectfont
		\textbf{QUESTION: Who is the most important proponent of deontology?}
		
		\vspace{0.1cm}
		\textbf{OPTIONS:}
		\begin{enumerate}
			\item Jeremy Bentham
			\item Immanuel Kant
			\item John Stuart Mill
			\item Aristotle
		\end{enumerate}
		
		\vspace{0.1cm}
		\textcolor{myblue}{\textbf{ANSWER: B}}
		
		\vspace{0.1cm}
		\textbf{DETAILED EXPLANATION:}
		\begin{itemize}
			\item \textbf{Why B is correct:} Immanuel Kant is the most important proponent of deontology. His categorical imperative states we should act only according to maxims that could be willed as universal law.
			\item \textbf{Why A is wrong:} Bentham is associated with utilitarianism
			\item \textbf{Why C is wrong:} Mill is associated with utilitarianism
			\item \textbf{Why D is wrong:} Aristotle is associated with virtue ethics
		\end{itemize}
		
		\textcolor{myred}{\textbf{EXAM TRAP:}} Key deontologist = \textcolor{myblue}{Immanuel Kant}!
	\end{frame}
	
	\begin{frame}{Q24: Categorical Imperative - Tutorial}
		\fontsize{6.5}{7.5}\selectfont
		\textbf{QUESTION: What is the categorical imperative in Kantian ethics?}
		
		\vspace{0.1cm}
		\textbf{OPTIONS:}
		\begin{enumerate}
			\item A rule that says consequences determine morality
			\item A fundamental principle that one should act only according to maxims that could be willed as universal law
			\item A principle that says virtue is most important
			\item A rule that says ethics is subjective
		\end{enumerate}
		
		\vspace{0.1cm}
		\textcolor{myblue}{\textbf{ANSWER: B}}
		
		\vspace{0.1cm}
		\textbf{DETAILED EXPLANATION:}
		\begin{itemize}
			\item \textbf{Why B is correct:} The categorical imperative asserts one should act only according to maxims that could be willed as a universal law without contradiction, such as "It's wrong to lie."
			\item \textbf{Why A is wrong:} This is consequentialism
			\item \textbf{Why C is wrong:} This is virtue ethics
			\item \textbf{Why D is wrong:} Kant believed ethics was objective
		\end{itemize}
		
		\textcolor{myred}{\textbf{EXAM TRAP:}} Categorical imperative = \textcolor{myblue}{act on maxims that could be universal law}!
	\end{frame}
	
	\begin{frame}{Q25: Treating People as Ends - Tutorial}
		\fontsize{6.5}{7.5}\selectfont
		\textbf{QUESTION: What does Kant mean by treating people as ends in themselves?}
		
		\vspace{0.1cm}
		\textbf{OPTIONS:}
		\begin{enumerate}
			\item Treating people as tools for your purposes
			\item Respecting that people have their own values and goals as autonomous beings
			\item Using people to achieve your goals
			\item Ignoring people's autonomy
		\end{enumerate}
		
		\vspace{0.1cm}
		\textcolor{myblue}{\textbf{ANSWER: B}}
		
		\vspace{0.1cm}
		\textbf{DETAILED EXPLANATION:}
		\begin{itemize}
			\item \textbf{Why B is correct:} Kant argued we shouldn't treat people like things. People are ends in themselves, having their own values and goals as autonomous beings.
			\item \textbf{Why A is wrong:} This is treating people as means, not ends
			\item \textbf{Why C is wrong:} This violates Kant's principle
			\item \textbf{Why D is wrong:} This disregards Kant's emphasis on autonomy
		\end{itemize}
		
		\textcolor{myred}{\textbf{EXAM TRAP:}} Treating people as ends = \textcolor{myblue}{respecting their autonomy and values}!
	\end{frame}
	
	\begin{frame}{Q26: Deontology Criticism - Tutorial}
		\fontsize{6.5}{7.5}\selectfont
		\textbf{QUESTION: What is a criticism of deontology?}
		
		\vspace{0.1cm}
		\textbf{OPTIONS:}
		\begin{enumerate}
			\item It focuses too much on consequences
			\item It can lead to rigid rule-following even when consequences are disastrous
			\item It ignores duties entirely
			\item It is too flexible
		\end{enumerate}
		
		\vspace{0.1cm}
		\textcolor{myblue}{\textbf{ANSWER: B}}
		
		\vspace{0.1cm}
		\textbf{DETAILED EXPLANATION:}
		\begin{itemize}
			\item \textbf{Why B is correct:} Deontology can lead to rigid rule-following even when consequences are disastrous. Kant famously argued lying is always wrong regardless of consequences.
			\item \textbf{Why A is wrong:} Deontology focuses on rules, not consequences
			\item \textbf{Why C is wrong:} Duties are central to deontology
			\item \textbf{Why D is wrong:} Deontology is often criticized as too rigid, not too flexible
		\end{itemize}
		
		\textcolor{myred}{\textbf{EXAM TRAP:}} Deontology can be \textcolor{myblue}{too rigid} even when consequences are disastrous!
	\end{frame}
	
	\begin{frame}{Q27: Virtue Ethics Definition - Tutorial}
		\fontsize{6.5}{7.5}\selectfont
		\textbf{QUESTION: What is virtue ethics?}
		
		\vspace{0.1cm}
		\textbf{OPTIONS:}
		\begin{enumerate}
			\item An ethical theory that judges morality based on consequences
			\item An ethical framework that judges morality based on adherence to rules
			\item An ethical theory that emphasizes virtuous character traits and living a good life
			\item An ethical theory that ignores character
		\end{enumerate}
		
		\vspace{0.1cm}
		\textcolor{myblue}{\textbf{ANSWER: C}}
		
		\vspace{0.1cm}
		\textbf{DETAILED EXPLANATION:}
		\begin{itemize}
			\item \textbf{Why C is correct:} Virtue ethics emphasizes virtuous character traits and living a good life, rather than focusing on rules or consequences. It asks "What kind of person should I be?"
			\item \textbf{Why A is wrong:} This is consequentialism
			\item \textbf{Why B is wrong:} This is deontology
			\item \textbf{Why D is wrong:} Character is central to virtue ethics
		\end{itemize}
		
		\textcolor{myred}{\textbf{EXAM TRAP:}} Virtue ethics = \textcolor{myblue}{focus on character and living a good life}!
	\end{frame}
	
	\begin{frame}{Q28: Aristotle and Virtue Ethics - Tutorial}
		\fontsize{6.5}{7.5}\selectfont
		\textbf{QUESTION: Who is the ancient Greek philosopher associated with virtue ethics?}
		
		\vspace{0.1cm}
		\textbf{OPTIONS:}
		\begin{enumerate}
			\item Plato
			\item Aristotle
			\item Socrates
			\item Epicurus
		\end{enumerate}
		
		\vspace{0.1cm}
		\textcolor{myblue}{\textbf{ANSWER: B}}
		
		\vspace{0.1cm}
		\textbf{DETAILED EXPLANATION:}
		\begin{itemize}
			\item \textbf{Why B is correct:} Virtue ethics has roots in ancient Greek philosophers like Aristotle, who taught that happiness comes from living a life guided by virtues.
			\item \textbf{Why A is wrong:} Plato is associated with different philosophical traditions
			\item \textbf{Why C is wrong:} Socrates is associated with different philosophical traditions
			\item \textbf{Why D is wrong:} Epicurus is associated with different philosophical traditions
		\end{itemize}
		
		\textcolor{myred}{\textbf{EXAM TRAP:}} Key virtue ethicist = \textcolor{myblue}{Aristotle}!
	\end{frame}
	
	\begin{frame}{Q29: Virtues Examples - Tutorial}
		\fontsize{6.5}{7.5}\selectfont
		\textbf{QUESTION: Which is an example of a virtue in virtue ethics?}
		
		\vspace{0.1cm}
		\textbf{OPTIONS:}
		\begin{enumerate}
			\item Maximizing utility
			\item Following all rules
			\item Wisdom, courage, humanity, justice, temperance, and generosity
			\item Calculating consequences
		\end{enumerate}
		
		\vspace{0.1cm}
		\textcolor{myblue}{\textbf{ANSWER: C}}
		
		\vspace{0.1cm}
		\textbf{DETAILED EXPLANATION:}
		\begin{itemize}
			\item \textbf{Why C is correct:} Virtues include wisdom, courage, humanity, justice, temperance, and generosity. These are character traits we should cultivate to live well.
			\item \textbf{Why A is wrong:} This is consequentialism
			\item \textbf{Why B is wrong:} This is deontology
			\item \textbf{Why D is wrong:} This is consequentialist thinking
		\end{itemize}
		
		\textcolor{myred}{\textbf{EXAM TRAP:}} Virtues include \textcolor{myblue}{wisdom, courage, justice, temperance, and generosity}!
	\end{frame}
	
	\begin{frame}{Q30: Virtue Ethics Key Question - Tutorial}
		\fontsize{6.5}{7.5}\selectfont
		\textbf{QUESTION: What is the key question in virtue ethics?}
		
		\vspace{0.1cm}
		\textbf{OPTIONS:}
		\begin{enumerate}
			\item "What action maximizes utility?"
			\item "What duty must I fulfill?"
			\item "What kind of person should I be?"
			\item "What are the consequences?"
		\end{enumerate}
		
		\vspace{0.1cm}
		\textcolor{myblue}{\textbf{ANSWER: C}}
		
		\vspace{0.1cm}
		\textbf{DETAILED EXPLANATION:}
		\begin{itemize}
			\item \textbf{Why C is correct:} Virtue ethics asks "What kind of person should I be?" rather than focusing on universal duties or maximizing utility.
			\item \textbf{Why A is wrong:} This is the consequentialist question
			\item \textbf{Why B is wrong:} This is the deontological question
			\item \textbf{Why D is wrong:} This is consequentialist thinking
		\end{itemize}
		
		\textcolor{myred}{\textbf{EXAM TRAP:}} Virtue ethics asks \textcolor{myblue}{"What kind of person should I be?"}!
	\end{frame}
	
	\begin{frame}{Q31: Virtue Ethics Criticism - Tutorial}
		\fontsize{6.5}{7.5}\selectfont
		\textbf{QUESTION: What is a criticism of virtue ethics?}
		
		\vspace{0.1cm}
		\textbf{OPTIONS:}
		\begin{enumerate}
			\item It provides too much clear guidance
			\item It lacks clear guidance for moral decisions compared to duty-based or consequentialist approaches
			\item It is too focused on consequences
			\item It is too focused on rules
		\end{enumerate}
		
		\vspace{0.1cm}
		\textcolor{myblue}{\textbf{ANSWER: B}}
		
		\vspace{0.1cm}
		\textbf{DETAILED EXPLANATION:}
		\begin{itemize}
			\item \textbf{Why B is correct:} Critics argue virtue ethics lacks clear guidance for moral decisions compared to duty-based or consequentialist approaches. Different virtues can conflict, and how to weigh them is unclear.
			\item \textbf{Why A is wrong:} It's often criticized for lack of clear guidance
			\item \textbf{Why C is wrong:} Virtue ethics focuses on character, not consequences
			\item \textbf{Why D is wrong:} Virtue ethics focuses on character, not rules
		\end{itemize}
		
		\textcolor{myred}{\textbf{EXAM TRAP:}} Virtue ethics criticized for \textcolor{myblue}{lack of clear guidance}!
	\end{frame}
	
	\begin{frame}{Q32: Virtue Ethics and Practical Wisdom - Tutorial}
		\fontsize{6.5}{7.5}\selectfont
		\textbf{QUESTION: How do virtue ethicists respond to criticism about lack of guidance?}
		
		\vspace{0.1cm}
		\textbf{OPTIONS:}
		\begin{enumerate}
			\item By ignoring the criticism
			\item By arguing that practical wisdom helps navigate hard cases
			\item By switching to consequentialism
			\item By adopting strict rules
		\end{enumerate}
		
		\vspace{0.1cm}
		\textcolor{myblue}{\textbf{ANSWER: B}}
		
		\vspace{0.1cm}
		\textbf{DETAILED EXPLANATION:}
		\begin{itemize}
			\item \textbf{Why B is correct:} Virtue ethicists counter that practical wisdom helps navigate hard cases. Moral duties can also conflict, and consequences are not always comparable.
			\item \textbf{Why A is wrong:} They engage with the criticism
			\item \textbf{Why C is wrong:} They don't abandon virtue ethics
			\item \textbf{Why D is wrong:} This would be adopting a deontological approach
		\end{itemize}
		
		\textcolor{myred}{\textbf{EXAM TRAP:}} Virtue ethics relies on \textcolor{myblue}{practical wisdom} to navigate hard cases!
	\end{frame}
	
	\begin{frame}{Q33: Modern Virtue Ethics Applications - Tutorial}
		\fontsize{6.5}{7.5}\selectfont
		\textbf{QUESTION: How has virtue ethics been extended to organizations?}
		
		\vspace{0.1cm}
		\textbf{OPTIONS:}
		\begin{enumerate}
			\item It has not been extended
			\item Organizational virtues might include justice, accountability, environmental stewardship, and responsible innovation
			\item Organizations cannot have virtues
			\item Virtue ethics only applies to individuals
		\end{enumerate}
		
		\vspace{0.1cm}
		\textcolor{myblue}{\textbf{ANSWER: B}}
		
		\vspace{0.1cm}
		\textbf{DETAILED EXPLANATION:}
		\begin{itemize}
			\item \textbf{Why B is correct:} Modern virtue ethics extends beyond individual character. For organizations, virtues might include justice, accountability, environmental stewardship, and responsible innovation.
			\item \textbf{Why A is wrong:} Virtue ethics has been extended to organizations
			\item \textbf{Why C is wrong:} Organizations can embody virtues through their practices
			\item \textbf{Why D is wrong:} Virtue ethics applies to organizations too
		\end{itemize}
		
		\textcolor{myred}{\textbf{EXAM TRAP:}} Organizational virtues include \textcolor{myblue}{justice, accountability, and responsible innovation}!
	\end{frame}
	
	\begin{frame}{Q34: Virtue Ethics and AI - Tutorial}
		\fontsize{6.5}{7.5}\selectfont
		\textbf{QUESTION: What is one way virtue ethics applies to AI development?}
		
		\vspace{0.1cm}
		\textbf{OPTIONS:}
		\begin{enumerate}
			\item By programming all possible rules into AI
			\item By having AI learn from experience and habit, like children do
			\item By maximizing utility through AI
			\item By ignoring character in AI
		\end{enumerate}
		
		\vspace{0.1cm}
		\textcolor{myblue}{\textbf{ANSWER: B}}
		
		\vspace{0.1cm}
		\textbf{DETAILED EXPLANATION:}
		\begin{itemize}
			\item \textbf{Why B is correct:} Some scholars suggest that to build ethical AI, we must build it in a virtue ethics way, which would imply it learning from experience and habit, like children do.
			\item \textbf{Why A is wrong:} This is a rule-based approach, not virtue ethics
			\item \textbf{Why C is wrong:} This is consequentialism
			\item \textbf{Why D is wrong:} Character is central to virtue ethics
		\end{itemize}
		
		\textcolor{myred}{\textbf{EXAM TRAP:}} Virtue ethics AI would \textcolor{myblue}{learn from experience like children}!
	\end{frame}
	
	\begin{frame}{Q35: Ethical Frameworks Complementarity - Tutorial}
		\fontsize{6.5}{7.5}\selectfont
		\textbf{QUESTION: How do consequentialism, deontology, and virtue ethics relate?}
		
		\vspace{0.1cm}
		\textbf{OPTIONS:}
		\begin{enumerate}
			\item They are mutually exclusive and cannot be combined
			\item They are not mutually exclusive and complement one another in practical ethics
			\item Only one framework is ever useful
			\item They all produce identical recommendations
		\end{enumerate}
		
		\vspace{0.1cm}
		\textcolor{myblue}{\textbf{ANSWER: B}}
		
		\vspace{0.1cm}
		\textbf{DETAILED EXPLANATION:}
		\begin{itemize}
			\item \textbf{Why B is correct:} Consequentialism, deontology, and virtue ethics are not mutually exclusive. In practical ethics, they complement one another. The best moral decision accords with all three theories.
			\item \textbf{Why A is wrong:} They can be combined
			\item \textbf{Why C is wrong:} Each provides valuable insights
			\item \textbf{Why D is wrong:} They have different emphases
		\end{itemize}
		
		\textcolor{myred}{\textbf{EXAM TRAP:}} Ethical frameworks are \textcolor{myblue}{complementary, not mutually exclusive}!
	\end{frame}
	
	% ============================================================
	% SECTION 4: MEDICAL ETHICS AND AI - QUESTIONS 36-40
	% ============================================================
	
	\begin{frame}{Q36: Medical Ethics History - Tutorial}
		\fontsize{6.5}{7.5}\selectfont
		\textbf{QUESTION: What is the Hippocratic Oath and its significance?}
		
		\vspace{0.1cm}
		\textbf{OPTIONS:}
		\begin{enumerate}
			\item A modern AI ethics document
			\item An ancient Greek oath emphasizing physicians doing no harm to patients
			\item A legal document from the 20th century
			\item A business ethics code
		\end{enumerate}
		
		\vspace{0.1cm}
		\textcolor{myblue}{\textbf{ANSWER: B}}
		
		\vspace{0.1cm}
		\textbf{DETAILED EXPLANATION:}
		\begin{itemize}
			\item \textbf{Why B is correct:} The Hippocratic Oath, written between the fifth and third centuries BC, emphasizes the importance of physicians doing no harm to patients. It's one of the earliest medical ethics documents.
			\item \textbf{Why A is wrong:} It's ancient Greek, not modern AI
			\item \textbf{Why C is wrong:} It's from ancient Greece
			\item \textbf{Why D is wrong:} It's medical, not business ethics
		\end{itemize}
		
		\textcolor{myred}{\textbf{EXAM TRAP:}} Hippocratic Oath = \textcolor{myblue}{"do no harm"} in medicine!
	\end{frame}
	
	\begin{frame}{Q37: Medical Ethics Development - Tutorial}
		\fontsize{6.5}{7.5}\selectfont
		\textbf{QUESTION: What factors drove the development of medical ethics in the 1970s?}
		
		\vspace{0.1cm}
		\textbf{OPTIONS:}
		\begin{enumerate}
			\item Only technological advances
			\item Concentration of medical care, rising costs, patient rights movements, public outrage at scandals, and rapid technological advances
			\item Only government regulation
			\item Only public opinion
		\end{enumerate}
		
		\vspace{0.1cm}
		\textcolor{myblue}{\textbf{ANSWER: B}}
		
		\vspace{0.1cm}
		\textbf{DETAILED EXPLANATION:}
		\begin{itemize}
			\item \textbf{Why B is correct:} Medical ethics evolved in the 1970s due to: concentration of medical care, rising costs, patient rights movements, public outrage at scandals (like Tuskegee), and rapid technological advances.
			\item \textbf{Why A is wrong:} Multiple factors contributed
			\item \textbf{Why C is wrong:} More than just regulation
			\item \textbf{Why D is wrong:} Multiple factors beyond public opinion
		\end{itemize}
		
		\textcolor{myred}{\textbf{EXAM TRAP:}} Medical ethics evolved due to \textcolor{myblue}{multiple societal and technological factors}!
	\end{frame}
	
	\begin{frame}{Q38: Medical Ethics Documents - Tutorial}
		\fontsize{6.5}{7.5}\selectfont
		\textbf{QUESTION: Which is NOT a major medical ethics document?}
		
		\vspace{0.1cm}
		\textbf{OPTIONS:}
		\begin{enumerate}
			\item Nuremberg Code (1947)
			\item Declaration of Helsinki (1964)
			\item Belmont Report (1978)
			\item EU AI Act (2024)
		\end{enumerate}
		
		\vspace{0.1cm}
		\textcolor{myblue}{\textbf{ANSWER: D}}
		
		\vspace{0.1cm}
		\textbf{DETAILED EXPLANATION:}
		\begin{itemize}
			\item \textbf{Why D is correct:} The EU AI Act is an AI regulation, not a medical ethics document. The others are all major medical ethics documents.
			\item \textbf{Why A is wrong:} Nuremberg Code IS a medical ethics document
			\item \textbf{Why B is wrong:} Declaration of Helsinki IS a medical ethics document
			\item \textbf{Why C is wrong:} Belmont Report IS a medical ethics document
		\end{itemize}
		
		\textcolor{myred}{\textbf{EXAM TRAP:}} EU AI Act is \textcolor{myblue}{AI regulation}, not medical ethics!
	\end{frame}
	
	\begin{frame}{Q39: Medical Ethics and AI Parallels - Tutorial}
		\fontsize{6.5}{7.5}\selectfont
		\textbf{QUESTION: Why is medical ethics relevant to AI ethics?}
		
		\vspace{0.1cm}
		\textbf{OPTIONS:}
		\begin{enumerate}
			\item They are completely unrelated
			\item Medical ethics has a longer history of addressing practical ethical dilemmas in high-stakes contexts, providing lessons for AI
			\item Medicine and AI use the same technologies
			\item Doctors are the only ones who use AI
		\end{enumerate}
		
		\vspace{0.1cm}
		\textcolor{myblue}{\textbf{ANSWER: B}}
		
		\vspace{0.1cm}
		\textbf{DETAILED EXPLANATION:}
		\begin{itemize}
			\item \textbf{Why B is correct:} Medical ethics has a longer history of addressing practical ethical dilemmas in high-stakes contexts (life and death). AI ethics can learn from these established frameworks and structures.
			\item \textbf{Why A is wrong:} They share important parallels
			\item \textbf{Why C is wrong:} They use different technologies
			\item \textbf{Why D is wrong:} AI is used across many fields
		\end{itemize}
		
		\textcolor{myred}{\textbf{EXAM TRAP:}} Medical ethics provides \textcolor{myblue}{valuable lessons} for AI ethics!
	\end{frame}
	
	\begin{frame}{Q40: Medical Ethics Lessons for AI - Tutorial}
		\fontsize{6.5}{7.5}\selectfont
		\textbf{QUESTION: What lesson can AI ethics learn from medical ethics?}
		
		\vspace{0.1cm}
		\textbf{OPTIONS:}
		\begin{enumerate}
			\item That ethics is unnecessary
			\item The importance of established ethical frameworks, structures, and committees to address complex dilemmas
			\item That technology always solves ethical problems
			\item That only doctors should make ethical decisions
		\end{enumerate}
		
		\vspace{0.1cm}
		\textcolor{myblue}{\textbf{ANSWER: B}}
		
		\vspace{0.1cm}
		\textbf{DETAILED EXPLANATION:}
		\begin{itemize}
			\item \textbf{Why B is correct:} Medical ethics shows the importance of established ethical frameworks, structures (ethics committees), and codes of ethics to address complex dilemmas - lessons AI ethics is still learning.
			\item \textbf{Why A is wrong:} Ethics is essential
			\item \textbf{Why C is wrong:} Technology doesn't solve ethical problems
			\item \textbf{Why D is wrong:} Ethical decisions involve many stakeholders
		\end{itemize}
		
		\textcolor{myred}{\textbf{EXAM TRAP:}} AI ethics needs \textcolor{myblue}{established structures and frameworks} like medical ethics!
	\end{frame}
	
	% ============================================================
	% SECTION 5: PRINCIPLES OF AI ETHICS - QUESTIONS 41-55
	% ============================================================
	
	\begin{frame}{Q41: AI Ethics Principles - Tutorial}
		\fontsize{6.5}{7.5}\selectfont
		\textbf{QUESTION: What are five common principles in AI ethics codes?}
		
		\vspace{0.1cm}
		\textbf{OPTIONS:}
		\begin{enumerate}
			\item Speed, efficiency, profit, growth, and market share
			\item Nonmaleficence, beneficence, justice, autonomy, and explainability
			\item Privacy, security, transparency, accountability, and fairness
			\item All of the above are common principles
		\end{enumerate}
		
		\vspace{0.1cm}
		\textcolor{myblue}{\textbf{ANSWER: B}}
		
		\vspace{0.1cm}
		\textbf{DETAILED EXPLANATION:}
		\begin{itemize}
			\item \textbf{Why B is correct:} The five principles common to many AI ethics codes are nonmaleficence (do no harm), beneficence (do good), justice (fairness), autonomy (respect for persons), and explainability (transparency).
			\item \textbf{Why A is wrong:} These are business metrics, not ethics principles
			\item \textbf{Why C is wrong:} These are related but not the five common principles
			\item \textbf{Why D is wrong:} The specific five are nonmaleficence, beneficence, justice, autonomy, and explainability
		\end{itemize}
		
		\textcolor{myred}{\textbf{EXAM TRAP:}} Five principles = \textcolor{myblue}{nonmaleficence, beneficence, justice, autonomy, explainability}!
	\end{frame}
	
	\begin{frame}{Q42: Nonmaleficence Definition - Tutorial}
		\fontsize{6.5}{7.5}\selectfont
		\textbf{QUESTION: What is the principle of nonmaleficence?}
		
		\vspace{0.1cm}
		\textbf{OPTIONS:}
		\begin{enumerate}
			\item The obligation to act for the benefit of others
			\item The obligation to avoid harming others or inflicting injuries
			\item The obligation to treat people fairly
			\item The obligation to respect people's autonomy
		\end{enumerate}
		
		\vspace{0.1cm}
		\textcolor{myblue}{\textbf{ANSWER: B}}
		
		\vspace{0.1cm}
		\textbf{DETAILED EXPLANATION:}
		\begin{itemize}
			\item \textbf{Why B is correct:} Nonmaleficence asserts an obligation to avoid harming others or inflicting injuries, including not imposing unjustified risks of harm.
			\item \textbf{Why A is wrong:} This is beneficence
			\item \textbf{Why C is wrong:} This is justice
			\item \textbf{Why D is wrong:} This is autonomy
		\end{itemize}
		
		\textcolor{myred}{\textbf{EXAM TRAP:}} Nonmaleficence = \textcolor{myblue}{"do no harm"} - avoid harming others!
	\end{frame}
	
	\begin{frame}{Q43: Beneficence Definition - Tutorial}
		\fontsize{6.5}{7.5}\selectfont
		\textbf{QUESTION: What is the principle of beneficence?}
		
		\vspace{0.1cm}
		\textbf{OPTIONS:}
		\begin{enumerate}
			\item The obligation to avoid harming others
			\item The moral obligation to act for the benefit of others
			\item The obligation to treat people fairly
			\item The obligation to respect people's autonomy
		\end{enumerate}
		
		\vspace{0.1cm}
		\textcolor{myblue}{\textbf{ANSWER: B}}
		
		\vspace{0.1cm}
		\textbf{DETAILED EXPLANATION:}
		\begin{itemize}
			\item \textbf{Why B is correct:} Beneficence refers to the moral obligation to act for the benefit of others. It requires taking positive steps to help others, not just refraining from harm.
			\item \textbf{Why A is wrong:} This is nonmaleficence
			\item \textbf{Why C is wrong:} This is justice
			\item \textbf{Why D is wrong:} This is autonomy
		\end{itemize}
		
		\textcolor{myred}{\textbf{EXAM TRAP:}} Beneficence = \textcolor{myblue}{"do good"} - act for the benefit of others!
	\end{frame}
	
	\begin{frame}{Q44: Nonmaleficence vs Beneficence - Tutorial}
		\fontsize{6.5}{7.5}\selectfont
		\textbf{QUESTION: How do nonmaleficence and beneficence differ?}
		
		\vspace{0.1cm}
		\textbf{OPTIONS:}
		\begin{enumerate}
			\item They are the same principle
			\item Nonmaleficence means avoiding harm; beneficence means actively doing good
			\item Nonmaleficence is more important than beneficence
			\item Beneficence is a subset of nonmaleficence
		\end{enumerate}
		
		\vspace{0.1cm}
		\textcolor{myblue}{\textbf{ANSWER: B}}
		
		\vspace{0.1cm}
		\textbf{DETAILED EXPLANATION:}
		\begin{itemize}
			\item \textbf{Why B is correct:} Nonmaleficence is the duty to avoid harming others (do no harm). Beneficence is the duty to actively do good and help others (do good). Beneficence goes beyond mere avoidance of harm.
			\item \textbf{Why A is wrong:} They are distinct principles
			\item \textbf{Why C is wrong:} Both are important in different ways
			\item \textbf{Why D is wrong:} Beneficence is broader than just avoiding harm
		\end{itemize}
		
		\textcolor{myred}{\textbf{EXAM TRAP:}} Nonmaleficence = \textcolor{myblue}{avoid harm}; Beneficence = \textcolor{myblue}{do good}!
	\end{frame}
	
	\begin{frame}{Q45: Justice Definition - Tutorial}
		\fontsize{6.5}{7.5}\selectfont
		\textbf{QUESTION: What is the principle of justice in ethics?}
		
		\vspace{0.1cm}
		\textbf{OPTIONS:}
		\begin{enumerate}
			\item The obligation to avoid harming others
			\item The obligation to act for the benefit of others
			\item The moral obligation to act in accordance with principles of fairness, equality, impartiality, and proportionality
			\item The obligation to respect people's autonomy
		\end{enumerate}
		
		\vspace{0.1cm}
		\textcolor{myblue}{\textbf{ANSWER: C}}
		
		\vspace{0.1cm}
		\textbf{DETAILED EXPLANATION:}
		\begin{itemize}
			\item \textbf{Why C is correct:} Justice refers to the moral obligation to act in accordance with principles of fairness, equality, impartiality, and proportionality. It requires giving each person their proper due.
			\item \textbf{Why A is wrong:} This is nonmaleficence
			\item \textbf{Why B is wrong:} This is beneficence
			\item \textbf{Why D is wrong:} This is autonomy
		\end{itemize}
		
		\textcolor{myred}{\textbf{EXAM TRAP:}} Justice = \textcolor{myblue}{fairness, equality, impartiality, and proportionality}!
	\end{frame}
	
	\begin{frame}{Q46: Types of Justice - Tutorial}
		\fontsize{6.5}{7.5}\selectfont
		\textbf{QUESTION: What is distributive justice?}
		
		\vspace{0.1cm}
		\textbf{OPTIONS:}
		\begin{enumerate}
			\item Fair processes and impartiality
			\item Equitable allocation of benefits and burdens in society
			\item Repairing harms through reconciling victims and offenders
			\item Respect and fairness between individuals
		\end{enumerate}
		
		\vspace{0.1cm}
		\textcolor{myblue}{\textbf{ANSWER: B}}
		
		\vspace{0.1cm}
		\textbf{DETAILED EXPLANATION:}
		\begin{itemize}
			\item \textbf{Why B is correct:} Distributive justice focuses on equitable allocation of benefits and burdens in society. It concerns how resources and opportunities are distributed.
			\item \textbf{Why A is wrong:} This is procedural justice
			\item \textbf{Why C is wrong:} This is restorative justice
			\item \textbf{Why D is wrong:} This is interactional justice
		\end{itemize}
		
		\textcolor{myred}{\textbf{EXAM TRAP:}} Distributive justice = \textcolor{myblue}{equitable allocation of benefits and burdens}!
	\end{frame}
	
	\begin{frame}{Q47: Procedural Justice - Tutorial}
		\fontsize{6.5}{7.5}\selectfont
		\textbf{QUESTION: What is procedural justice?}
		
		\vspace{0.1cm}
		\textbf{OPTIONS:}
		\begin{enumerate}
			\item Equitable allocation of benefits and burdens
			\item Fair processes and impartiality
			\item Repairing harms through reconciliation
			\item Respect and fairness between individuals
		\end{enumerate}
		
		\vspace{0.1cm}
		\textcolor{myblue}{\textbf{ANSWER: B}}
		
		\vspace{0.1cm}
		\textbf{DETAILED EXPLANATION:}
		\begin{itemize}
			\item \textbf{Why B is correct:} Procedural justice demands fair processes and impartiality. It's about ensuring that the procedures used to make decisions are fair, not just the outcomes.
			\item \textbf{Why A is wrong:} This is distributive justice
			\item \textbf{Why C is wrong:} This is restorative justice
			\item \textbf{Why D is wrong:} This is interactional justice
		\end{itemize}
		
		\textcolor{myred}{\textbf{EXAM TRAP:}} Procedural justice = \textcolor{myblue}{fair processes and impartiality}!
	\end{frame}
	
	\begin{frame}{Q48: Autonomy Definition - Tutorial}
		\fontsize{6.5}{7.5}\selectfont
		\textbf{QUESTION: What is the principle of autonomy in ethics?}
		
		\vspace{0.1cm}
		\textbf{OPTIONS:}
		\begin{enumerate}
			\item The capacity to make informed, un-coerced decisions about one's life
			\item The obligation to avoid harming others
			\item The obligation to act for the benefit of others
			\item The obligation to treat people fairly
		\end{enumerate}
		
		\vspace{0.1cm}
		\textcolor{myblue}{\textbf{ANSWER: A}}
		
		\vspace{0.1cm}
		\textbf{DETAILED EXPLANATION:}
		\begin{itemize}
			\item \textbf{Why A is correct:} Autonomy refers to the capacity of people to make their own informed, un-coerced decisions about their lives and actions. It commands respecting others' abilities to determine their own course.
			\item \textbf{Why B is wrong:} This is nonmaleficence
			\item \textbf{Why C is wrong:} This is beneficence
			\item \textbf{Why D is wrong:} This is justice
		\end{itemize}
		
		\textcolor{myred}{\textbf{EXAM TRAP:}} Autonomy = \textcolor{myblue}{making informed, un-coerced decisions}!
	\end{frame}
	
	\begin{frame}{Q49: Autonomy in Healthcare - Tutorial}
		\fontsize{6.5}{7.5}\selectfont
		\textbf{QUESTION: How is autonomy applied in healthcare?}
		
		\vspace{0.1cm}
		\textbf{OPTIONS:}
		\begin{enumerate}
			\item Patients must always follow doctor's orders
			\item Patients have a right to voluntary informed consent and refusal regarding their treatment
			\item Patients have no say in their treatment
			\item Doctors make all decisions for patients
		\end{enumerate}
		
		\vspace{0.1cm}
		\textcolor{myblue}{\textbf{ANSWER: B}}
		
		\vspace{0.1cm}
		\textbf{DETAILED EXPLANATION:}
		\begin{itemize}
			\item \textbf{Why B is correct:} In healthcare, respect for autonomy is crucial. Patients have a right to voluntary informed consent and refusal regarding their treatment. Doctors must inform, but patients decide.
			\item \textbf{Why A is wrong:} Patients have autonomy rights
			\item \textbf{Why C is wrong:} Patients have a say in their treatment
			\item \textbf{Why D is wrong:} Shared decision-making is preferred
		\end{itemize}
		
		\textcolor{myred}{\textbf{EXAM TRAP:}} Autonomy in healthcare = \textcolor{myblue}{patient consent and refusal rights}!
	\end{frame}
	
	\begin{frame}{Q50: Autonomy in AI - Tutorial}
		\fontsize{6.5}{7.5}\selectfont
		\textbf{QUESTION: How does autonomy apply to ethical AI?}
		
		\vspace{0.1cm}
		\textbf{OPTIONS:}
		\begin{enumerate}
			\item AI should make all decisions for people
			\item Ethical AI respects people's autonomy by not using coercive or manipulative tactics
			\item Autonomy is irrelevant to AI
			\item AI should override human decisions
		\end{enumerate}
		
		\vspace{0.1cm}
		\textcolor{myblue}{\textbf{ANSWER: B}}
		
		\vspace{0.1cm}
		\textbf{DETAILED EXPLANATION:}
		\begin{itemize}
			\item \textbf{Why B is correct:} An ethical AI system respects people's autonomy by not using coercive or manipulative tactics to get people to act in a particular way. Technology should help people further their own goals.
			\item \textbf{Why A is wrong:} This would violate autonomy
			\item \textbf{Why C is wrong:} Autonomy is relevant to AI
			\item \textbf{Why D is wrong:} This would violate autonomy
		\end{itemize}
		
		\textcolor{myred}{\textbf{EXAM TRAP:}} Ethical AI respects \textcolor{myblue}{autonomy by avoiding coercion and manipulation}!
	\end{frame}
	
	\begin{frame}{Q51: Explainability Definition - Tutorial}
		\fontsize{6.5}{7.5}\selectfont
		\textbf{QUESTION: What is the principle of explainability in AI?}
		
		\vspace{0.1cm}
		\textbf{OPTIONS:}
		\begin{enumerate}
			\item Making AI systems faster
			\item AI systems should provide transparent and understandable explanations for their actions or decisions
			\item Making AI systems more accurate
			\item Making AI systems cheaper
		\end{enumerate}
		
		\vspace{0.1cm}
		\textcolor{myblue}{\textbf{ANSWER: B}}
		
		\vspace{0.1cm}
		\textbf{DETAILED EXPLANATION:}
		\begin{itemize}
			\item \textbf{Why B is correct:} Explainability (or explicability) refers to the idea that AI systems, especially those with decision-making capabilities, should provide transparent and understandable explanations for their actions or decisions.
			\item \textbf{Why A is wrong:} Speed is not explainability
			\item \textbf{Why C is wrong:} Accuracy is not explainability
			\item \textbf{Why D is wrong:} Cost is not explainability
		\end{itemize}
		
		\textcolor{myred}{\textbf{EXAM TRAP:}} Explainability = \textcolor{myblue}{transparent and understandable explanations}!
	\end{frame}
	
	\begin{frame}{Q52: Importance of Explainability - Tutorial}
		\fontsize{6.5}{7.5}\selectfont
		\textbf{QUESTION: Why is explainability important in AI?}
		
		\vspace{0.1cm}
		\textbf{OPTIONS:}
		\begin{enumerate}
			\item It makes AI more profitable
			\item It enables accountability, builds trust, and helps identify biases
			\item It makes AI faster
			\item It reduces costs
		\end{enumerate}
		
		\vspace{0.1cm}
		\textcolor{myblue}{\textbf{ANSWER: B}}
		
		\vspace{0.1cm}
		\textbf{DETAILED EXPLANATION:}
		\begin{itemize}
			\item \textbf{Why B is correct:} Explainability is important for accountability (holding designers accountable), building trust (users understand decisions), and identifying biases and errors.
			\item \textbf{Why A is wrong:} Profit is not the primary reason
			\item \textbf{Why C is wrong:} Speed is not the primary reason
			\item \textbf{Why D is wrong:} Cost reduction is not the primary reason
		\end{itemize}
		
		\textcolor{myred}{\textbf{EXAM TRAP:}} Explainability enables \textcolor{myblue}{accountability, trust, and bias detection}!
	\end{frame}
	
	\begin{frame}{Q53: Types of Explainability - Tutorial}
		\fontsize{6.5}{7.5}\selectfont
		\textbf{QUESTION: What is local explainability in AI?}
		
		\vspace{0.1cm}
		\textbf{OPTIONS:}
		\begin{enumerate}
			\item Explaining the overall behavior of a model
			\item Explaining a specific decision for a single instance or prediction
			\item Explaining the architecture of a model
			\item Explaining the training data
		\end{enumerate}
		
		\vspace{0.1cm}
		\textcolor{myblue}{\textbf{ANSWER: B}}
		
		\vspace{0.1cm}
		\textbf{DETAILED EXPLANATION:}
		\begin{itemize}
			\item \textbf{Why B is correct:} Local explainability focuses on explaining the decisions of a specific AI model on a single instance or prediction. It provides insights into why a particular decision was made.
			\item \textbf{Why A is wrong:} This is global explainability
			\item \textbf{Why C is wrong:} This is model architecture explanation
			\item \textbf{Why D is wrong:} This is data explanation
		\end{itemize}
		
		\textcolor{myred}{\textbf{EXAM TRAP:}} Local explainability = \textcolor{myblue}{explaining a single decision/prediction}!
	\end{frame}
	
	\begin{frame}{Q54: Global Explainability - Tutorial}
		\fontsize{6.5}{7.5}\selectfont
		\textbf{QUESTION: What is global explainability in AI?}
		
		\vspace{0.1cm}
		\textbf{OPTIONS:}
		\begin{enumerate}
			\item Explaining a specific decision for a single instance
			\item Looking at an AI model's overall behavior and decision-making processes
			\item Explaining the training data
			\item Explaining the model's architecture
		\end{enumerate}
		
		\vspace{0.1cm}
		\textcolor{myblue}{\textbf{ANSWER: B}}
		
		\vspace{0.1cm}
		\textbf{DETAILED EXPLANATION:}
		\begin{itemize}
			\item \textbf{Why B is correct:} Global explainability looks at an AI model's overall behavior and decision-making processes, providing a comprehensive understanding of how the model operates across various inputs.
			\item \textbf{Why A is wrong:} This is local explainability
			\item \textbf{Why C is wrong:} This is data explanation
			\item \textbf{Why D is wrong:} This is architecture explanation
		\end{itemize}
		
		\textcolor{myred}{\textbf{EXAM TRAP:}} Global explainability = \textcolor{myblue}{overall model behavior and processes}!
	\end{frame}
	
	\begin{frame}{Q55: Counterfactual Explanations - Tutorial}
		\fontsize{6.5}{7.5}\selectfont
		\textbf{QUESTION: What is a counterfactual explanation in AI?}
		
		\vspace{0.1cm}
		\textbf{OPTIONS:}
		\begin{enumerate}
			\item An explanation that ignores hypotheticals
			\item A hypothetical scenario that contrasts with the actual decision, showing what would have led to a different outcome
			\item An explanation that only describes the actual decision
			\item An explanation that blames the user
		\end{enumerate}
		
		\vspace{0.1cm}
		\textcolor{myblue}{\textbf{ANSWER: B}}
		
		\vspace{0.1cm}
		\textbf{DETAILED EXPLANATION:}
		\begin{itemize}
			\item \textbf{Why B is correct:} A counterfactual explanation presents a hypothetical scenario that contrasts with the actual decision. For example, "You would have been approved if you had \$10,000 more in the bank."
			\item \textbf{Why A is wrong:} Counterfactuals are based on hypotheticals
			\item \textbf{Why C is wrong:} It includes alternatives, not just actual decisions
			\item \textbf{Why D is wrong:} It provides actionable insights, not blame
		\end{itemize}
		
		\textcolor{myred}{\textbf{EXAM TRAP:}} Counterfactual explanations show \textcolor{myblue}{what would have led to a different outcome}!
	\end{frame}
	
	% ============================================================
	% SECTION 6: BIAS, DISCRIMINATION, AND FAIRNESS - QUESTIONS 56-75
	% ============================================================
	
	\begin{frame}{Q56: Bias in AI - Tutorial}
		\fontsize{6.5}{7.5}\selectfont
		\textbf{QUESTION: What is bias in the context of AI?}
		
		\vspace{0.1cm}
		\textbf{OPTIONS:}
		\begin{enumerate}
			\item Random errors in algorithmic predictions
			\item A systemic deviation in the output or impact of an algorithm compared to a desired norm or standard
			\item When algorithms are always perfectly accurate
			\item When algorithms are faster than humans
		\end{enumerate}
		
		\vspace{0.1cm}
		\textcolor{myblue}{\textbf{ANSWER: B}}
		
		\vspace{0.1cm}
		\textbf{DETAILED EXPLANATION:}
		\begin{itemize}
			\item \textbf{Why B is correct:} Bias in AI refers to a systemic deviation in the output or impact of an algorithm compared to a desired norm or standard. An algorithm is biased when it deviates from its intended function.
			\item \textbf{Why A is wrong:} Bias is systemic, not random
			\item \textbf{Why C is wrong:} Algorithms can be biased
			\item \textbf{Why D is wrong:} Speed doesn't determine bias
		\end{itemize}
		
		\textcolor{myred}{\textbf{EXAM TRAP:}} AI bias = \textcolor{myblue}{systemic deviation from desired norm}!
	\end{frame}
	
	\begin{frame}{Q57: Justifiable Bias - Tutorial}
		\fontsize{6.5}{7.5}\selectfont
		\textbf{QUESTION: Are all biases in AI problematic?}
		
		\vspace{0.1cm}
		\textbf{OPTIONS:}
		\begin{enumerate}
			\item Yes, all biases are problematic
			\item No, not all biases are problematic; justifiable biases can form part of a well-designed AI
			\item Bias never occurs in AI
			\item Bias is always intentional
		\end{enumerate}
		
		\vspace{0.1cm}
		\textcolor{myblue}{\textbf{ANSWER: B}}
		
		\vspace{0.1cm}
		\textbf{DETAILED EXPLANATION:}
		\begin{itemize}
			\item \textbf{Why B is correct:} Not all biases are problematic from an ethical standpoint. Justifiable biases can form part of a well-designed AI (e.g., prioritizing certain relevant criteria).
			\item \textbf{Why A is wrong:} Some biases are justifiable
			\item \textbf{Why C is wrong:} Bias does occur in AI
			\item \textbf{Why D is wrong:} Bias is often unintentional
		\end{itemize}
		
		\textcolor{myred}{\textbf{EXAM TRAP:}} Not all biases are \textcolor{myblue}{problematic} - some are justifiable!
	\end{frame}
	
	\begin{frame}{Q58: Problem Specification Bias - Tutorial}
		\fontsize{6.5}{7.5}\selectfont
		\textbf{QUESTION: How can bias arise in problem specification?}
		
		\vspace{0.1cm}
		\textbf{OPTIONS:}
		\begin{enumerate}
			\item Problems are always specified perfectly
			\item Operationalizing complex goals is nuanced, and selected target variables may fail to capture real-world objectives accurately
			\item Problem specification never introduces bias
			\item All specifications are equally valid
		\end{enumerate}
		
		\vspace{0.1cm}
		\textcolor{myblue}{\textbf{ANSWER: B}}
		
		\vspace{0.1cm}
		\textbf{DETAILED EXPLANATION:}
		\begin{itemize}
			\item \textbf{Why B is correct:} An algorithm may exhibit bias from its inception if the goals it is designed to achieve contain inherent problems. Selected target variables may fail to capture real-world objectives accurately.
			\item \textbf{Why A is wrong:} Problems are often specified imperfectly
			\item \textbf{Why C is wrong:} Problem specification CAN introduce bias
			\item \textbf{Why D is wrong:} Specifications vary in quality
		\end{itemize}
		
		\textcolor{myred}{\textbf{EXAM TRAP:}} Problem specification can introduce \textcolor{myblue}{bias through poor operationalization}!
	\end{frame}
	
	\begin{frame}{Q59: Healthcare Expenditure Proxy Example - Tutorial}
		\fontsize{6.5}{7.5}\selectfont
		\textbf{QUESTION: What happened with the algorithm using healthcare expenditure as a proxy?}
		
		\vspace{0.1cm}
		\textbf{OPTIONS:}
		\begin{enumerate}
			\item It perfectly identified the sickest patients
			\item It favored the richest patients who spend more on healthcare, not the sickest ones
			\item It was completely unbiased
			\item It was never deployed
		\end{enumerate}
		
		\vspace{0.1cm}
		\textcolor{myblue}{\textbf{ANSWER: B}}
		
		\vspace{0.1cm}
		\textbf{DETAILED EXPLANATION:}
		\begin{itemize}
			\item \textbf{Why B is correct:} An algorithm meant to identify sicker patients used healthcare expenditure as a proxy, thereby favoring not the sickest patients, but the richest ones who tend to spend more on healthcare.
			\item \textbf{Why A is wrong:} It failed to identify the sickest patients
			\item \textbf{Why C is wrong:} It was biased
			\item \textbf{Why D is wrong:} The example shows it was deployed
		\end{itemize}
		
		\textcolor{myred}{\textbf{EXAM TRAP:}} Proxies can \textcolor{myblue}{introduce bias} by favoring different groups than intended!
	\end{frame}
	
	\begin{frame}{Q60: Historical Bias Definition - Tutorial}
		\fontsize{6.5}{7.5}\selectfont
		\textbf{QUESTION: What is historical bias in AI?}
		
		\vspace{0.1cm}
		\textbf{OPTIONS:}
		\begin{enumerate}
			\item Bias that occurs only in the present
			\item Bias in data that reflects past discrimination and patterns
			\item Bias that is always intentional
			\item Bias that never affects algorithms
		\end{enumerate}
		
		\vspace{0.1cm}
		\textcolor{myblue}{\textbf{ANSWER: B}}
		
		\vspace{0.1cm}
		\textbf{DETAILED EXPLANATION:}
		\begin{itemize}
			\item \textbf{Why B is correct:} Historical bias occurs when ML algorithms rely on historical data that reflects past discrimination and biases (e.g., women historically excluded from banking).
			\item \textbf{Why A is wrong:} Historical bias is about the PAST
			\item \textbf{Why C is wrong:} Historical bias is often unintentional
			\item \textbf{Why D is wrong:} Historical bias DOES affect algorithms
		\end{itemize}
		
		\textcolor{myred}{\textbf{EXAM TRAP:}} Historical bias = \textcolor{myblue}{past discrimination reflected in data}!
	\end{frame}
	
	\begin{frame}{Q61: Selective Labels Problem - Tutorial}
		\fontsize{6.5}{7.5}\selectfont
		\textbf{QUESTION: What is the selective labels problem in AI?}
		
		\vspace{0.1cm}
		\textbf{OPTIONS:}
		\begin{enumerate}
			\item All data has complete labels
			\item Historical data rarely shows counterfactual outcomes (e.g., what would have happened if we had hired or lent to different people)
			\item Labels are always accurate
			\item Labels are never needed
		\end{enumerate}
		
		\vspace{0.1cm}
		\textcolor{myblue}{\textbf{ANSWER: B}}
		
		\vspace{0.1cm}
		\textbf{DETAILED EXPLANATION:}
		\begin{itemize}
			\item \textbf{Why B is correct:} Historical data rarely shows counterfactual outcomes. A company knows about people it hired or lent to, but not about those it rejected, creating a selective labels problem.
			\item \textbf{Why A is wrong:} Data often has incomplete labels
			\item \textbf{Why C is wrong:} Labels can be inaccurate
			\item \textbf{Why D is wrong:} Labels are needed for supervised learning
		\end{itemize}
		
		\textcolor{myred}{\textbf{EXAM TRAP:}} Selective labels = \textcolor{myblue}{missing counterfactual outcomes} (what would have happened)!
	\end{frame}
	
	\begin{frame}{Q62: Sampling Bias - Tutorial}
		\fontsize{6.5}{7.5}\selectfont
		\textbf{QUESTION: What is sampling bias in AI?}
		
		\vspace{0.1cm}
		\textbf{OPTIONS:}
		\begin{enumerate}
			\item When data samples are perfectly random
			\item When the data sample is not random, so trends may not generalize to other populations
			\item When all groups are equally represented
			\item When data is always representative
		\end{enumerate}
		
		\vspace{0.1cm}
		\textcolor{myblue}{\textbf{ANSWER: B}}
		
		\vspace{0.1cm}
		\textbf{DETAILED EXPLANATION:}
		\begin{itemize}
			\item \textbf{Why B is correct:} Sampling bias arises when the data sample is not random. If the sample is not random, the trends shown by the population under study may not generalize to another population.
			\item \textbf{Why A is wrong:} Sampling bias means NOT random
			\item \textbf{Why C is wrong:} This would be balanced sampling
			\item \textbf{Why D is wrong:} Sampling bias means NOT representative
		\end{itemize}
		
		\textcolor{myred}{\textbf{EXAM TRAP:}} Sampling bias = \textcolor{myblue}{non-random sample that doesn't generalize}!
	\end{frame}
	
	\begin{frame}{Q63: Modeling and Design Bias - Tutorial}
		\fontsize{6.5}{7.5}\selectfont
		\textbf{QUESTION: How can bias emerge in modeling and design?}
		
		\vspace{0.1cm}
		\textbf{OPTIONS:}
		\begin{enumerate}
			\item Modeling never introduces bias
			\item Choices related to optimization functions, regression models, subgroup consideration, and information presentation can all introduce biases
			\item Only data introduces bias
			\item Only deployment introduces bias
		\end{enumerate}
		
		\vspace{0.1cm}
		\textcolor{myblue}{\textbf{ANSWER: B}}
		
		\vspace{0.1cm}
		\textbf{DETAILED EXPLANATION:}
		\begin{itemize}
			\item \textbf{Why B is correct:} Even with sound problem specification and unbiased data, biases can emerge in modeling, validation, and design phases through choices of optimization functions, models, subgroup consideration, and information presentation.
			\item \textbf{Why A is wrong:} Modeling CAN introduce bias
			\item \textbf{Why C is wrong:} Bias can emerge at multiple stages
			\item \textbf{Why D is wrong:} Bias can emerge at multiple stages
		\end{itemize}
		
		\textcolor{myred}{\textbf{EXAM TRAP:}} Modeling choices can \textcolor{myblue}{introduce bias} even with good data!
	\end{frame}
	
	\begin{frame}{Q64: Deployment Bias - Tutorial}
		\fontsize{6.5}{7.5}\selectfont
		\textbf{QUESTION: How can bias emerge during deployment?}
		
		\vspace{0.1cm}
		\textbf{OPTIONS:}
		\begin{enumerate}
			\item Deployment never introduces bias
			\item People may defer entirely to algorithm suggestions, creating unintended effects
			\item Algorithms always work as intended in deployment
			\item Deployment eliminates all bias
		\end{enumerate}
		
		\vspace{0.1cm}
		\textcolor{myblue}{\textbf{ANSWER: B}}
		
		\vspace{0.1cm}
		\textbf{DETAILED EXPLANATION:}
		\begin{itemize}
			\item \textbf{Why B is correct:} Even when all other aspects are addressed, biases can emerge when people defer entirely to algorithm suggestions, creating unintended effects and allowing blame-shifting.
			\item \textbf{Why A is wrong:} Deployment CAN introduce bias
			\item \textbf{Why C is wrong:} Algorithms may not work as intended
			\item \textbf{Why D is wrong:} Deployment doesn't eliminate bias
		\end{itemize}
		
		\textcolor{myred}{\textbf{EXAM TRAP:}} Deployment can introduce \textcolor{myblue}{bias through human deference to algorithms}!
	\end{frame}
	
	\begin{frame}{Q65: Automation Bias in Decision Making - Tutorial}
		\fontsize{6.5}{7.5}\selectfont
		\textbf{QUESTION: Why do people tend to defer to automated systems?}
		
		\vspace{0.1cm}
		\textbf{OPTIONS:}
		\begin{enumerate}
			\item Because they are always correct
			\item Because it's convenient, algorithms appear more objective, and it may shield people from responsibility
			\item Because they are forced to
			\item Because algorithms are never wrong
		\end{enumerate}
		
		\vspace{0.1cm}
		\textcolor{myblue}{\textbf{ANSWER: B}}
		
		\vspace{0.1cm}
		\textbf{DETAILED EXPLANATION:}
		\begin{itemize}
			\item \textbf{Why B is correct:} People defer to automated systems because it's convenient and time-saving, algorithms appear more "objective," and deferring may shield people from responsibility if things go wrong.
			\item \textbf{Why A is wrong:} Algorithms are not always correct
			\item \textbf{Why C is wrong:} They are not forced to
			\item \textbf{Why D is wrong:} Algorithms can be wrong
		\end{itemize}
		
		\textcolor{myred}{\textbf{EXAM TRAP:}} People defer to algorithms because of \textcolor{myblue}{convenience, perceived objectivity, and responsibility shielding}!
	\end{frame}
	
	\begin{frame}{Q66: Discrimination Legal Perspective - Tutorial}
		\fontsize{6.5}{7.5}\selectfont
		\textbf{QUESTION: When does algorithmic bias count as discrimination?}
		
		\vspace{0.1cm}
		\textbf{OPTIONS:}
		\begin{enumerate}
			\item Always
			\item Never
			\item When it results in disfavoring people based on protected characteristics (race, sex, ethnicity, age, etc.) in violation of law
			\item Only when intentional
		\end{enumerate}
		
		\vspace{0.1cm}
		\textcolor{myblue}{\textbf{ANSWER: C}}
		
		\vspace{0.1cm}
		\textbf{DETAILED EXPLANATION:}
		\begin{itemize}
			\item \textbf{Why C is correct:} Algorithmic bias is likely to lead to discrimination when it results in disfavoring people based on protected characteristics (race, sex, ethnicity, age, etc.) in violation of legal protections.
			\item \textbf{Why A is wrong:} Not all bias is discriminatory
			\item \textbf{Why B is wrong:} Some bias IS discriminatory
			\item \textbf{Why D is wrong:} Discrimination can be unintentional
		\end{itemize}
		
		\textcolor{myred}{\textbf{EXAM TRAP:}} Discrimination = \textcolor{myblue}{disfavoring based on protected characteristics} in violation of law!
	\end{frame}
	
	\begin{frame}{Q67: Group Fairness - Tutorial}
		\fontsize{6.5}{7.5}\selectfont
		\textbf{QUESTION: What is group fairness?}
		
		\vspace{0.1cm}
		\textbf{OPTIONS:}
		\begin{enumerate}
			\item Treating similar individuals similarly
			\item Statistical criteria where, for example, loan demographics mirror the overall population
			\item Treating every individual uniquely
			\item Ignoring group-level patterns
		\end{enumerate}
		
		\vspace{0.1cm}
		\textcolor{myblue}{\textbf{ANSWER: B}}
		
		\vspace{0.1cm}
		\textbf{DETAILED EXPLANATION:}
		\begin{itemize}
			\item \textbf{Why B is correct:} Group fairness involves statistical criteria, such as requiring that the demographics of approved individuals mirror the overall population (e.g., if 51% are female, 51% of loan recipients should be female).
			\item \textbf{Why A is wrong:} This is individual fairness
			\item \textbf{Why C is wrong:} This is not fairness
			\item \textbf{Why D is wrong:} Group fairness examines group patterns
		\end{itemize}
		
		\textcolor{myred}{\textbf{EXAM TRAP:}} Group fairness = \textcolor{myblue}{statistical equality across groups}!
	\end{frame}
	
	\begin{frame}{Q68: Individual Fairness - Tutorial}
		\fontsize{6.5}{7.5}\selectfont
		\textbf{QUESTION: What is individual fairness?}
		
		\vspace{0.1cm}
		\textbf{OPTIONS:}
		\begin{enumerate}
			\item Statistical criteria across groups
			\item Treating similar individuals similarly
			\item Ignoring individual differences
			\item Treating everyone the same regardless of context
		\end{enumerate}
		
		\vspace{0.1cm}
		\textcolor{myblue}{\textbf{ANSWER: B}}
		
		\vspace{0.1cm}
		\textbf{DETAILED EXPLANATION:}
		\begin{itemize}
			\item \textbf{Why B is correct:} Individual fairness emphasizes treating similar individuals similarly, even though defining similarity can pose challenges.
			\item \textbf{Why A is wrong:} This is group fairness
			\item \textbf{Why C is wrong:} Individual differences matter in fairness
			\item \textbf{Why D is wrong:} Context matters in fairness
		\end{itemize}
		
		\textcolor{myred}{\textbf{EXAM TRAP:}} Individual fairness = \textcolor{myblue}{treating similar individuals similarly}!
	\end{frame}
	
	\begin{frame}{Q69: Fairness Impossibility - Tutorial}
		\fontsize{6.5}{7.5}\selectfont
		\textbf{QUESTION: Why can't fairness be automated when base rates are unequal?}
		
		\vspace{0.1cm}
		\textbf{OPTIONS:}
		\begin{enumerate}
			\item Because algorithms are not powerful enough
			\item Because it's mathematically impossible to satisfy multiple operationalizations of fairness simultaneously when base rates differ
			\item Because fairness doesn't matter
			\item Because data is always perfect
		\end{enumerate}
		
		\vspace{0.1cm}
		\textcolor{myblue}{\textbf{ANSWER: B}}
		
		\vspace{0.1cm}
		\textbf{DETAILED EXPLANATION:}
		\begin{itemize}
			\item \textbf{Why B is correct:} When base rates between populations are different, it is mathematically impossible to satisfy demographic parity, predictive rate parity, and equal opportunity simultaneously.
			\item \textbf{Why A is wrong:} It's a mathematical impossibility, not a power issue
			\item \textbf{Why C is wrong:} Fairness does matter
			\item \textbf{Why D is wrong:} Data is never perfect
		\end{itemize}
		
		\textcolor{myred}{\textbf{EXAM TRAP:}} Fairness cannot be automated when \textcolor{myblue}{base rates are unequal} - mathematical impossibility!
	\end{frame}
	
	\begin{frame}{Q70: Procedural Fairness - Tutorial}
		\fontsize{6.5}{7.5}\selectfont
		\textbf{QUESTION: What is procedural fairness in AI?}
		
		\vspace{0.1cm}
		\textbf{OPTIONS:}
		\begin{enumerate}
			\item Ensuring fair outcomes only
			\item Ensuring fair processes and impartial procedures
			\item Ignoring processes entirely
			\item Focusing only on consequences
		\end{enumerate}
		
		\vspace{0.1cm}
		\textcolor{myblue}{\textbf{ANSWER: B}}
		
		\vspace{0.1cm}
		\textbf{DETAILED EXPLANATION:}
		\begin{itemize}
			\item \textbf{Why B is correct:} Procedural fairness provides reassurances that fair outcomes will be sought through impartial and just processes. It's about having the right structures in place.
			\item \textbf{Why A is wrong:} Both process and outcome matter
			\item \textbf{Why C is wrong:} Processes are important
			\item \textbf{Why D is wrong:} Both process and consequences matter
		\end{itemize}
		
		\textcolor{myred}{\textbf{EXAM TRAP:}} Procedural fairness = \textcolor{myblue}{fair processes and impartial procedures}!
	\end{frame}
	
	\begin{frame}{Q71: Avoiding Problematic Biases - Tutorial}
		\fontsize{6.5}{7.5}\selectfont
		\textbf{QUESTION: What is the goal when dealing with algorithmic biases?}
		
		\vspace{0.1cm}
		\textbf{OPTIONS:}
		\begin{enumerate}
			\item To avoid all biases completely
			\item To avoid problematic biases, recognizing that some biases are justifiable
			\item To ignore biases entirely
			\item To eliminate human judgment
		\end{enumerate}
		
		\vspace{0.1cm}
		\textcolor{myblue}{\textbf{ANSWER: B}}
		
		\vspace{0.1cm}
		\textbf{DETAILED EXPLANATION:}
		\begin{itemize}
			\item \textbf{Why B is correct:} Since there is no such thing as an "objective" algorithm and fairness cannot be fully automated, the task is not to avoid biases in general, but to avoid problematic biases.
			\item \textbf{Why A is wrong:} Avoiding all biases is impossible and undesirable
			\item \textbf{Why C is wrong:} Biases cannot be ignored
			\item \textbf{Why D is wrong:} Human judgment remains important
		\end{itemize}
		
		\textcolor{myred}{\textbf{EXAM TRAP:}} Goal = \textcolor{myblue}{avoid problematic biases} - some biases are justifiable!
	\end{frame}
	
	\begin{frame}{Q72: Technological Solutions for Bias - Tutorial}
		\fontsize{6.5}{7.5}\selectfont
		\textbf{QUESTION: Can technology alone solve bias problems?}
		
		\vspace{0.1cm}
		\textbf{OPTIONS:}
		\begin{enumerate}
			\item Yes, technology can solve all bias problems
			\item No, technology rarely solves bias problems on its own; fairness cannot be fully automated
			\item No technology exists for bias detection
			\item Bias is not a problem technology can address
		\end{enumerate}
		
		\vspace{0.1cm}
		\textcolor{myblue}{\textbf{ANSWER: B}}
		
		\vspace{0.1cm}
		\textbf{DETAILED EXPLANATION:}
		\begin{itemize}
			\item \textbf{Why B is correct:} Although technological tools can be helpful, technology rarely solves bias problems on its own. Fairness cannot be fully automated, and AI cannot solve the fairness problems it creates.
			\item \textbf{Why A is wrong:} Technology alone is insufficient
			\item \textbf{Why C is wrong:} Technology does exist for bias detection
			\item \textbf{Why D is wrong:} Technology can help, but not solve alone
		\end{itemize}
		
		\textcolor{myred}{\textbf{EXAM TRAP:}} Technology alone \textcolor{myblue}{cannot solve bias} - fairness cannot be fully automated!
	\end{frame}
	
	\begin{frame}{Q73: Auditing for Bias - Tutorial}
		\fontsize{6.5}{7.5}\selectfont
		\textbf{QUESTION: What is likely the best way to identify and correct biases?}
		
		\vspace{0.1cm}
		\textbf{OPTIONS:}
		\begin{enumerate}
			\item Ignoring them
			\item Algorithmic auditing by independent assessors
			\item Only using technological tools
			\item Removing all features
		\end{enumerate}
		
		\vspace{0.1cm}
		\textcolor{myblue}{\textbf{ANSWER: B}}
		
		\vspace{0.1cm}
		\textbf{DETAILED EXPLANATION:}
		\begin{itemize}
			\item \textbf{Why B is correct:} Algorithmic auditing is likely the best way to identify and correct biases. It includes technological tools and statistical analyses, but also a fresh and diverse look at systems.
			\item \textbf{Why A is wrong:} Ignoring biases is not a solution
			\item \textbf{Why C is wrong:} Technology alone is insufficient
			\item \textbf{Why D is wrong:} Removing features may remove relevant information
		\end{itemize}
		
		\textcolor{myred}{\textbf{EXAM TRAP:}} Algorithmic auditing is \textcolor{myblue}{best way to identify and correct biases}!
	\end{frame}
	
	\begin{frame}{Q74: Data Quality and Bias - Tutorial}
		\fontsize{6.5}{7.5}\selectfont
		\textbf{QUESTION: How can data quality help avoid biases?}
		
		\vspace{0.1cm}
		\textbf{OPTIONS:}
		\begin{enumerate}
			\item Data quality is irrelevant to bias
			\item Ensuring data is diverse, updated, accurate, representative, and free from past discriminatory tendencies helps avoid biases
			\item Only synthetic data can avoid bias
			\item Data quality cannot help with bias
		\end{enumerate}
		
		\vspace{0.1cm}
		\textcolor{myblue}{\textbf{ANSWER: B}}
		
		\vspace{0.1cm}
		\textbf{DETAILED EXPLANATION:}
		\begin{itemize}
			\item \textbf{Why B is correct:} Ensuring data is diverse, updated, accurate, representative, and free from past discriminatory tendencies goes a long way toward avoiding biases.
			\item \textbf{Why A is wrong:} Data quality is important
			\item \textbf{Why C is wrong:} Synthetic data is an option, not the only solution
			\item \textbf{Why D is wrong:} Data quality CAN help
		\end{itemize}
		
		\textcolor{myred}{\textbf{EXAM TRAP:}} Data quality helps \textcolor{myblue}{avoid biases} through diversity and representativeness!
	\end{frame}
	
	\begin{frame}{Q75: Ethical Committees - Tutorial}
		\fontsize{6.5}{7.5}\selectfont
		\textbf{QUESTION: How can ethical committees help with fairness?}
		
		\vspace{0.1cm}
		\textbf{OPTIONS:}
		\begin{enumerate}
			\item They are unnecessary
			\item They can discuss possible problematic biases and make decisions about trade-offs considering different ethical frameworks
			\item They eliminate all biases
			\item They replace all technical work
		\end{enumerate}
		
		\vspace{0.1cm}
		\textcolor{myblue}{\textbf{ANSWER: B}}
		
		\vspace{0.1cm}
		\textbf{DETAILED EXPLANATION:}
		\begin{itemize}
			\item \textbf{Why B is correct:} Ethical committees provide forums where possible problematic biases can be discussed, and decisions about trade-offs can be made considering consequentialist, deontological, and virtue ethics considerations.
			\item \textbf{Why A is wrong:} They are valuable
			\item \textbf{Why C is wrong:} They can't eliminate all biases
			\item \textbf{Why D is wrong:} They complement technical work
		\end{itemize}
		
		\textcolor{myred}{\textbf{EXAM TRAP:}} Ethical committees help \textcolor{myblue}{discuss and address biases} through diverse perspectives!
	\end{frame}
	
	% ============================================================
	% SECTION 7: PRIVACY AND CYBERSECURITY - QUESTIONS 76-85
	% ============================================================
	
	\begin{frame}{Q76: Privacy Definition - Tutorial}
		\fontsize{6.5}{7.5}\selectfont
		\textbf{QUESTION: What is privacy in the context of AI?}
		
		\vspace{0.1cm}
		\textbf{OPTIONS:}
		\begin{enumerate}
			\item The ability to control all data
			\item Having privacy to the extent that others don't have access to our personal information
			\item Having complete freedom from all observation
			\item The right to be forgotten
		\end{enumerate}
		
		\vspace{0.1cm}
		\textcolor{myblue}{\textbf{ANSWER: B}}
		
		\vspace{0.1cm}
		\textbf{DETAILED EXPLANATION:}
		\begin{itemize}
			\item \textbf{Why B is correct:} In the context of AI, someone has privacy with respect to some person or institution and in reference to some personal data point if that person or institution has no access to that personal data point.
			\item \textbf{Why A is wrong:} Privacy is about access, not control
			\item \textbf{Why C is wrong:} Some observation is inevitable
			\item \textbf{Why D is wrong:} Right to be forgotten is part of privacy, not the whole definition
		\end{itemize}
		
		\textcolor{myred}{\textbf{EXAM TRAP:}} Privacy = \textcolor{myblue}{lack of access to personal information}!
	\end{frame}
	
	\begin{frame}{Q77: Data Security Definition - Tutorial}
		\fontsize{6.5}{7.5}\selectfont
		\textbf{QUESTION: What is data security?}
		
		\vspace{0.1cm}
		\textbf{OPTIONS:}
		\begin{enumerate}
			\item Doing what is necessary to prevent unauthorized access to data
			\item Ensuring data is always accessible
			\item Making data public
			\item Deleting all data
		\end{enumerate}
		
		\vspace{0.1cm}
		\textcolor{myblue}{\textbf{ANSWER: A}}
		
		\vspace{0.1cm}
		\textbf{DETAILED EXPLANATION:}
		\begin{itemize}
			\item \textbf{Why A is correct:} Data security is doing what is necessary to prevent unauthorized access to data. It includes protecting from ransomware, theft, and corruption.
			\item \textbf{Why B is wrong:} Accessibility must be balanced with security
			\item \textbf{Why C is wrong:} Making data public violates privacy
			\item \textbf{Why D is wrong:} Deleting all data is not data security
		\end{itemize}
		
		\textcolor{myred}{\textbf{EXAM TRAP:}} Data security = \textcolor{myblue}{preventing unauthorized access} to data!
	\end{frame}
	
	\begin{frame}{Q78: Privacy as an Ethical Issue - Tutorial}
		\fontsize{6.5}{7.5}\selectfont
		\textbf{QUESTION: Why is privacy an ethical issue?}
		
		\vspace{0.1cm}
		\textbf{OPTIONS:}
		\begin{enumerate}
			\item Because it's always illegal
			\item Because lack of privacy can lead to wrongs, harms, and risks for individuals, institutions, and society
			\item Because privacy is never important
			\item Because it only affects corporations
		\end{enumerate}
		
		\vspace{0.1cm}
		\textcolor{myblue}{\textbf{ANSWER: B}}
		
		\vspace{0.1cm}
		\textbf{DETAILED EXPLANATION:}
		\begin{itemize}
			\item \textbf{Why B is correct:} Privacy is an ethical issue because the lack of it can lead to wrongs, harms, and risks for individuals (discrimination, blackmail), institutions (lawsuits, fines), and society (national security risks).
			\item \textbf{Why A is wrong:} Privacy violations may be legal but still unethical
			\item \textbf{Why C is wrong:} Privacy IS important
			\item \textbf{Why D is wrong:} Privacy affects everyone, not just corporations
		\end{itemize}
		
		\textcolor{myred}{\textbf{EXAM TRAP:}} Privacy protects against \textcolor{myblue}{wrongs, harms, and risks} to individuals, institutions, and society!
	\end{frame}
	
	\begin{frame}{Q79: Personal Data as Toxic Asset - Tutorial}
		\fontsize{6.5}{7.5}\selectfont
		\textbf{QUESTION: Why is personal data described as a "toxic asset"?}
		
		\vspace{0.1cm}
		\textbf{OPTIONS:}
		\begin{enumerate}
			\item Because it's always worthless
			\item It's cheap, easy to mine, and profitable, but exposes companies to risk of hacks, leaks, lawsuits, and expensive management
			\item Because it's always harmful
			\item Because it's never valuable
		\end{enumerate}
		
		\vspace{0.1cm}
		\textcolor{myblue}{\textbf{ANSWER: B}}
		
		\vspace{0.1cm}
		\textbf{DETAILED EXPLANATION:}
		\begin{itemize}
			\item \textbf{Why B is correct:} Personal data is described as a toxic but potentially highly valuable asset. It's cheap, easy to mine, and profitable, but it exposes the company to risk of hacks, leaks, lawsuits, and is expensive to manage.
			\item \textbf{Why A is wrong:} It can be valuable
			\item \textbf{Why C is wrong:} It's not always harmful
			\item \textbf{Why D is wrong:} It can be very valuable
		\end{itemize}
		
		\textcolor{myred}{\textbf{EXAM TRAP:}} Personal data is a \textcolor{myblue}{toxic asset} - valuable but risky!
	\end{frame}
	
	\begin{frame}{Q80: Data Minimization - Tutorial}
		\fontsize{6.5}{7.5}\selectfont
		\textbf{QUESTION: What is data minimization?}
		
		\vspace{0.1cm}
		\textbf{OPTIONS:}
		\begin{enumerate}
			\item Collecting as much data as possible
			\item Collecting only the data that is necessary to fulfill a specific purpose
			\item Collecting no data at all
			\item Collecting data indefinitely
		\end{enumerate}
		
		\vspace{0.1cm}
		\textcolor{myblue}{\textbf{ANSWER: B}}
		
		\vspace{0.1cm}
		\textbf{DETAILED EXPLANATION:}
		\begin{itemize}
			\item \textbf{Why B is correct:} Data minimization is the most effective way to protect privacy: collecting only the data that is necessary to fulfill a specific purpose.
			\item \textbf{Why A is wrong:} This increases risk
			\item \textbf{Why C is wrong:} Some data is needed for services
			\item \textbf{Why D is wrong:} Data should be deleted when no longer needed
		\end{itemize}
		
		\textcolor{myred}{\textbf{EXAM TRAP:}} Data minimization = \textcolor{myblue}{collect only necessary data}!
	\end{frame}
	
	\begin{frame}{Q81: Right to be Forgotten - Tutorial}
		\fontsize{6.5}{7.5}\selectfont
		\textbf{QUESTION: What is the right to be forgotten?}
		
		\vspace{0.1cm}
		\textbf{OPTIONS:}
		\begin{enumerate}
			\item The right to have all data permanently stored
			\item A person's right to have private information removed from Internet searches when data is inadequate, irrelevant, or excessive
			\item The right to delete all company data
			\item The right to never be remembered
		\end{enumerate}
		
		\vspace{0.1cm}
		\textcolor{myblue}{\textbf{ANSWER: B}}
		
		\vspace{0.1cm}
		\textbf{DETAILED EXPLANATION:}
		\begin{itemize}
			\item \textbf{Why B is correct:} The right to be forgotten is a person's right to have private information about them removed from Internet searches and other directories when the data is inadequate, irrelevant, or excessive.
			\item \textbf{Why A is wrong:} This is the opposite
			\item \textbf{Why C is wrong:} It applies to personal information
			\item \textbf{Why D is wrong:} It's about removing information, not being completely forgotten
		\end{itemize}
		
		\textcolor{myred}{\textbf{EXAM TRAP:}} Right to be forgotten = \textcolor{myblue}{removal of private information when inadequate or irrelevant}!
	\end{frame}
	
	\begin{frame}{Q82: Contextual Integrity - Tutorial}
		\fontsize{6.5}{7.5}\selectfont
		\textbf{QUESTION: What is contextual integrity in privacy?}
		
		\vspace{0.1cm}
		\textbf{OPTIONS:}
		\begin{enumerate}
			\item Data can be used in any context
			\item The use of personal data should adhere to contextual norms of privacy
			\item Privacy norms are irrelevant
			\item Data should always be shared
		\end{enumerate}
		
		\vspace{0.1cm}
		\textcolor{myblue}{\textbf{ANSWER: B}}
		
		\vspace{0.1cm}
		\textbf{DETAILED EXPLANATION:}
		\begin{itemize}
			\item \textbf{Why B is correct:} When people give up their data in a particular context, they have certain expectations about how that data will be used. When transferred to a different context, privacy norms are violated.
			\item \textbf{Why A is wrong:} This violates contextual integrity
			\item \textbf{Why C is wrong:} Privacy norms are important
			\item \textbf{Why D is wrong:} Sharing should respect context
		\end{itemize}
		
		\textcolor{myred}{\textbf{EXAM TRAP:}} Contextual integrity = \textcolor{myblue}{data use should respect context-specific norms}!
	\end{frame}
	
	\begin{frame}{Q83: Data Deletion - Tutorial}
		\fontsize{6.5}{7.5}\selectfont
		\textbf{QUESTION: When should personal data be deleted?}
		
		\vspace{0.1cm}
		\textbf{OPTIONS:}
		\begin{enumerate}
			\item Never
			\item As soon as it is no longer necessary
			\item After 50 years
			\item Only when requested
		\end{enumerate}
		
		\vspace{0.1cm}
		\textcolor{myblue}{\textbf{ANSWER: B}}
		
		\vspace{0.1cm}
		\textbf{DETAILED EXPLANATION:}
		\begin{itemize}
			\item \textbf{Why B is correct:} Personal data should be deleted as soon as it is not necessary. Routine data deletion protects individuals and keeps data accurate (personal data changes quickly).
			\item \textbf{Why A is wrong:} Keeping data forever creates unnecessary risk
			\item \textbf{Why C is wrong:} 50 years is arbitrary and too long
			\item \textbf{Why D is wrong:} Proactive deletion is better than waiting for requests
		\end{itemize}
		
		\textcolor{myred}{\textbf{EXAM TRAP:}} Delete data \textcolor{myblue}{as soon as no longer necessary}!
	\end{frame}
	
	\begin{frame}{Q84: Data Security Practices - Tutorial}
		\fontsize{6.5}{7.5}\selectfont
		\textbf{QUESTION: What is a good data security practice?}
		
		\vspace{0.1cm}
		\textbf{OPTIONS:}
		\begin{enumerate}
			\item Using weak passwords
			\item Using strong encryption, thorough anonymization, and differential privacy
			\item Sharing data widely
			\item Storing data unencrypted
		\end{enumerate}
		
		\vspace{0.1cm}
		\textcolor{myblue}{\textbf{ANSWER: B}}
		
		\vspace{0.1cm}
		\textbf{DETAILED EXPLANATION:}
		\begin{itemize}
			\item \textbf{Why B is correct:} Good data security practices include using all technical tools available, from strong encryption and passwords to thorough anonymization and differential privacy.
			\item \textbf{Why A is wrong:} Weak passwords compromise security
			\item \textbf{Why C is wrong:} Sharing data widely increases risk
			\item \textbf{Why D is wrong:} Unencrypted data is vulnerable
		\end{itemize}
		
		\textcolor{myred}{\textbf{EXAM TRAP:}} Data security requires \textcolor{myblue}{encryption, anonymization, and differential privacy}!
	\end{frame}
	
	\begin{frame}{Q85: Privacy Principles Summary - Tutorial}
		\fontsize{6.5}{7.5}\selectfont
		\textbf{QUESTION: Which is NOT a key privacy principle?}
		
		\vspace{0.1cm}
		\textbf{OPTIONS:}
		\begin{enumerate}
			\item Right to privacy
			\item Data minimization
			\item Collecting as much data as possible
			\item Control over data
		\end{enumerate}
		
		\vspace{0.1cm}
		\textcolor{myblue}{\textbf{ANSWER: C}}
		
		\vspace{0.1cm}
		\textbf{DETAILED EXPLANATION:}
		\begin{itemize}
			\item \textbf{Why C is correct:} Collecting as much data as possible is the opposite of data minimization, which is a key privacy principle. It increases risk and violates privacy best practices.
			\item \textbf{Why A is wrong:} Right to privacy IS a key principle
			\item \textbf{Why B is wrong:} Data minimization IS a key principle
			\item \textbf{Why D is wrong:} Control over data IS a key principle
		\end{itemize}
		
		\textcolor{myred}{\textbf{EXAM TRAP:}} Collecting as much data as possible \textcolor{myblue}{violates privacy principles}!
	\end{frame}
	
	% ============================================================
	% SECTION 8: GOVERNANCE CHALLENGES - QUESTIONS 86-95
	% ============================================================
	
	\begin{frame}{Q86: AI Governance Challenges - Tutorial}
		\fontsize{6.5}{7.5}\selectfont
		\textbf{QUESTION: What makes AI difficult to govern?}
		
		\vspace{0.1cm}
		\textbf{OPTIONS:}
		\begin{enumerate}
			\item AI is always simple to understand
			\item Power asymmetries, lack of transparency, lack of interpretability, lack of ethics structures, lack of regulation, unpredictability, and privacy/copyright concerns
			\item AI never needs governance
			\item Only technical issues matter
		\end{enumerate}
		
		\vspace{0.1cm}
		\textcolor{myblue}{\textbf{ANSWER: B}}
		
		\vspace{0.1cm}
		\textbf{DETAILED EXPLANATION:}
		\begin{itemize}
			\item \textbf{Why B is correct:} AI governance is challenged by power asymmetries, lack of transparency, lack of interpretability, lack of AI ethics structures, lack of national and international regulation, unpredictability, lack of truth-tracking, and privacy/copyright concerns.
			\item \textbf{Why A is wrong:} AI is often complex and opaque
			\item \textbf{Why C is wrong:} AI requires governance
			\item \textbf{Why D is wrong:} Multiple governance challenges exist beyond technical issues
		\end{itemize}
		
		\textcolor{myred}{\textbf{EXAM TRAP:}} AI governance faces \textcolor{myblue}{multiple interconnected challenges}!
	\end{frame}
	
	\begin{frame}{Q87: Power Asymmetries - Tutorial}
		\fontsize{6.5}{7.5}\selectfont
		\textbf{QUESTION: What is a power asymmetry in AI governance?}
		
		\vspace{0.1cm}
		\textbf{OPTIONS:}
		\begin{enumerate}
			\item When governments are more powerful than tech companies
			\item Tech companies are often more powerful than some national governments and regulatory agencies
			\item When all stakeholders have equal power
			\item When power is never an issue
		\end{enumerate}
		
		\vspace{0.1cm}
		\textcolor{myblue}{\textbf{ANSWER: B}}
		
		\vspace{0.1cm}
		\textbf{DETAILED EXPLANATION:}
		\begin{itemize}
			\item \textbf{Why B is correct:} Many tech companies at the cutting edge of AI are often more powerful (in wealth, influence, and expertise) than some national governments and regulatory agencies, creating governance challenges.
			\item \textbf{Why A is wrong:} Often tech companies have more power
			\item \textbf{Why C is wrong:} Power is often unequal
			\item \textbf{Why D is wrong:} Power asymmetries ARE an issue
		\end{itemize}
		
		\textcolor{myred}{\textbf{EXAM TRAP:}} Tech companies may be \textcolor{myblue}{more powerful than some governments}!
	\end{frame}
	
	\begin{frame}{Q88: Lack of Transparency in AI - Tutorial}
		\fontsize{6.5}{7.5}\selectfont
		\textbf{QUESTION: Why is there a lack of transparency in AI systems?}
		
		\vspace{0.1cm}
		\textbf{OPTIONS:}
		\begin{enumerate}
			\item All AI systems are fully transparent
			\item Most AI systems have been developed by private companies not subject to the same transparency requirements as public institutions
			\item Transparency is never needed
			\item AI systems are always open-source
		\end{enumerate}
		
		\vspace{0.1cm}
		\textcolor{myblue}{\textbf{ANSWER: B}}
		
		\vspace{0.1cm}
		\textbf{DETAILED EXPLANATION:}
		\begin{itemize}
			\item \textbf{Why B is correct:} Most AI systems have been developed by private companies that may not be subject to the same transparency requirements as public institutions or universities.
			\item \textbf{Why A is wrong:} Most AI systems lack full transparency
			\item \textbf{Why C is wrong:} Transparency IS needed
			\item \textbf{Why D is wrong:} Most AI systems are proprietary
		\end{itemize}
		
		\textcolor{myred}{\textbf{EXAM TRAP:}} Private companies are \textcolor{myblue}{not subject to same transparency requirements} as public institutions!
	\end{frame}
	
	\begin{frame}{Q89: Lack of Interpretability - Tutorial}
		\fontsize{6.5}{7.5}\selectfont
		\textbf{QUESTION: Why are neural networks sometimes called "black boxes"?}
		
		\vspace{0.1cm}
		\textbf{OPTIONS:}
		\begin{enumerate}
			\item Because they are always transparent
			\item Because often not even computer scientists can be sure of exactly how the model is doing what it's doing
			\item Because they are simple to understand
			\item Because they never make mistakes
		\end{enumerate}
		
		\vspace{0.1cm}
		\textcolor{myblue}{\textbf{ANSWER: B}}
		
		\vspace{0.1cm}
		\textbf{DETAILED EXPLANATION:}
		\begin{itemize}
			\item \textbf{Why B is correct:} Neural networks are called "black boxes" because often not even computer scientists can be sure of exactly how the model is doing what it's doing due to their complexity.
			\item \textbf{Why A is wrong:} They are NOT transparent
			\item \textbf{Why C is wrong:} They are complex, not simple
			\item \textbf{Why D is wrong:} They can make mistakes
		\end{itemize}
		
		\textcolor{myred}{\textbf{EXAM TRAP:}} Neural networks are \textcolor{myblue}{black boxes} - even experts can't fully understand them!
	\end{frame}
	
	\begin{frame}{Q90: Proxy Variables and Interpretability - Tutorial}
		\fontsize{6.5}{7.5}\selectfont
		\textbf{QUESTION: How can proxy variables affect interpretability?}
		
		\vspace{0.1cm}
		\textbf{OPTIONS:}
		\begin{enumerate}
			\item They always make models more interpretable
			\item If outside observers don't know what proxies the model is using, it becomes difficult to govern the model
			\item Proxies never affect interpretability
			\item Proxies eliminate the black-box problem
		\end{enumerate}
		
		\vspace{0.1cm}
		\textcolor{myblue}{\textbf{ANSWER: B}}
		
		\vspace{0.1cm}
		\textbf{DETAILED EXPLANATION:}
		\begin{itemize}
			\item \textbf{Why B is correct:} The use of proxies can introduce issues, and if outside observers don't know what proxies the model is using, it is difficult to govern that model.
			\item \textbf{Why A is wrong:} Proxies can reduce interpretability
			\item \textbf{Why C is wrong:} Proxies DO affect interpretability
			\item \textbf{Why D is wrong:} Proxies don't eliminate the black-box problem
		\end{itemize}
		
		\textcolor{myred}{\textbf{EXAM TRAP:}} Unknown proxies \textcolor{myblue}{make governance difficult}!
	\end{frame}
	
	\begin{frame}{Q91: Lack of AI Ethics Structures - Tutorial}
		\fontsize{6.5}{7.5}\selectfont
		\textbf{QUESTION: How does AI ethics compare to medical ethics in terms of structure?}
		
		\vspace{0.1cm}
		\textbf{OPTIONS:}
		\begin{enumerate}
			\item AI ethics has more structure than medical ethics
			\item Medical ethics is supported by bioethicists, ethical codes, and committees; AI ethics lacks similar structures
			\item Both have identical structures
			\item AI ethics has no need for structures
		\end{enumerate}
		
		\vspace{0.1cm}
		\textcolor{myblue}{\textbf{ANSWER: B}}
		
		\vspace{0.1cm}
		\textbf{DETAILED EXPLANATION:}
		\begin{itemize}
			\item \textbf{Why B is correct:} Medical ethics is supported by bioethicists, ethical codes, ethics committees, licensing, and required coursework. AI ethics lacks similar structures - no required courses, no licensing, no committees.
			\item \textbf{Why A is wrong:} Medical ethics has MORE structure
			\item \textbf{Why C is wrong:} They are very different
			\item \textbf{Why D is wrong:} AI ethics needs structures
		\end{itemize}
		
		\textcolor{myred}{\textbf{EXAM TRAP:}} AI ethics lacks the \textcolor{myblue}{established structures} of medical ethics!
	\end{frame}
	
	\begin{frame}{Q92: Lack of Regulation - Tutorial}
		\fontsize{6.5}{7.5}\selectfont
		\textbf{QUESTION: What is the current state of AI regulation globally?}
		
		\vspace{0.1cm}
		\textbf{OPTIONS:}
		\begin{enumerate}
			\item Comprehensive global regulation exists
			\item There are still few national laws regulating AI (with the exception of China), and no specialized agencies overseeing AI
			\item All countries have comprehensive AI laws
			\item Regulation is unnecessary
		\end{enumerate}
		
		\vspace{0.1cm}
		\textcolor{myblue}{\textbf{ANSWER: B}}
		
		\vspace{0.1cm}
		\textbf{DETAILED EXPLANATION:}
		\begin{itemize}
			\item \textbf{Why B is correct:} At the time of writing, there are still no national laws regulating AI (with the exception of China). There are likewise no specialized agencies overseeing AI.
			\item \textbf{Why A is wrong:} Global regulation is limited
			\item \textbf{Why C is wrong:} Most countries lack AI laws
			\item \textbf{Why D is wrong:} Regulation is needed
		\end{itemize}
		
		\textcolor{myred}{\textbf{EXAM TRAP:}} Few national AI laws exist \textcolor{myblue}{(except China)} - no specialized agencies!
	\end{frame}
	
	\begin{frame}{Q93: Unpredictability in AI - Tutorial}
		\fontsize{6.5}{7.5}\selectfont
		\textbf{QUESTION: What is emergent behavior in AI systems?}
		
		\vspace{0.1cm}
		\textbf{OPTIONS:}
		\begin{enumerate}
			\item Behavior that is always predictable
			\item Behavior and properties that were not explicitly coded into the system
			\item Behavior that never surprises experts
			\item Behavior that is always fully understood
		\end{enumerate}
		
		\vspace{0.1cm}
		\textcolor{myblue}{\textbf{ANSWER: B}}
		
		\vspace{0.1cm}
		\textbf{DETAILED EXPLANATION:}
		\begin{itemize}
			\item \textbf{Why B is correct:} AI systems like neural networks can present emergent behavior and properties that were not explicitly coded into the system, surprising both laypeople and experts.
			\item \textbf{Why A is wrong:} Emergent behavior is UNPREDICTABLE
			\item \textbf{Why C is wrong:} It can surprise experts
			\item \textbf{Why D is wrong:} It may not be fully understood
		\end{itemize}
		
		\textcolor{myred}{\textbf{EXAM TRAP:}} Emergent behavior = \textcolor{myblue}{unexpected capabilities not explicitly coded}!
	\end{frame}
	
	\begin{frame}{Q94: Lack of Truth-Tracking - Tutorial}
		\fontsize{6.5}{7.5}\selectfont
		\textbf{QUESTION: Why are LLMs not truth-tracking?}
		
		\vspace{0.1cm}
		\textbf{OPTIONS:}
		\begin{enumerate}
			\item They are designed to always tell the truth
			\item They make statistical inferences and probabilistic "guesses" to construct plausible responses, not necessarily truthful ones
			\item They never make errors
			\item They are always accurate
		\end{enumerate}
		
		\vspace{0.1cm}
		\textcolor{myblue}{\textbf{ANSWER: B}}
		
		\vspace{0.1cm}
		\textbf{DETAILED EXPLANATION:}
		\begin{itemize}
			\item \textbf{Why B is correct:} LLMs are based on statistical inferences and probabilistic "guesses" to construct responses. They are designed to give plausible responses, but plausible responses are not necessarily truthful.
			\item \textbf{Why A is wrong:} LLMs are not designed for truth
			\item \textbf{Why C is wrong:} LLMs can make errors
			\item \textbf{Why D is wrong:} LLMs can be inaccurate
		\end{itemize}
		
		\textcolor{myred}{\textbf{EXAM TRAP:}} LLMs give \textcolor{myblue}{plausible, not necessarily truthful} responses!
	\end{frame}
	
	\begin{frame}{Q95: Privacy and Copyright Governance - Tutorial}
		\fontsize{6.5}{7.5}\selectfont
		\textbf{QUESTION: What are privacy and copyright concerns with AI?}
		
		\vspace{0.1cm}
		\textbf{OPTIONS:}
		\begin{enumerate}
			\item There are no concerns
			\item Companies scraped data off the internet without clear claims to that data, violating privacy and copyright
			\item All AI training data is public domain
			\item Privacy and copyright are not relevant to AI
		\end{enumerate}
		
		\vspace{0.1cm}
		\textcolor{myblue}{\textbf{ANSWER: B}}
		
		\vspace{0.1cm}
		\textbf{DETAILED EXPLANATION:}
		\begin{itemize}
			\item \textbf{Why B is correct:} When companies scrape data off the internet, questions arise about whether they have a claim to that data. Privacy violations and copyright infringement are major concerns.
			\item \textbf{Why A is wrong:} There are significant concerns
			\item \textbf{Why C is wrong:} AI training data is not public domain
			\item \textbf{Why D is wrong:} Privacy and copyright ARE relevant
		\end{itemize}
		
		\textcolor{myred}{\textbf{EXAM TRAP:}} Scraping data raises \textcolor{myblue}{privacy and copyright concerns}!
	\end{frame}
	
	% ============================================================
	% SECTION 9: REGULATORY LANDSCAPE - QUESTIONS 96-100
	% ============================================================
	
	\begin{frame}{Q96: GDPR Overview - Tutorial}
		\fontsize{6.5}{7.5}\selectfont
		\textbf{QUESTION: What is the GDPR?}
		
		\vspace{0.1cm}
		\textbf{OPTIONS:}
		\begin{enumerate}
			\item A US privacy law
			\item The European General Data Protection Regulation, designed to regulate personal data
			\item An AI-specific regulation
			\item A cybersecurity framework
		\end{enumerate}
		
		\vspace{0.1cm}
		\textcolor{myblue}{\textbf{ANSWER: B}}
		
		\vspace{0.1cm}
		\textbf{DETAILED EXPLANATION:}
		\begin{itemize}
			\item \textbf{Why B is correct:} The GDPR (General Data Protection Regulation) is a European regulation implemented in May 2018, designed to regulate personal data.
			\item \textbf{Why A is wrong:} It's European, not US
			\item \textbf{Why C is wrong:} It's about privacy, not AI specifically
			\item \textbf{Why D is wrong:} It's a privacy regulation, not just cybersecurity
		\end{itemize}
		
		\textcolor{myred}{\textbf{EXAM TRAP:}} GDPR = \textcolor{myblue}{European General Data Protection Regulation}!
	\end{frame}
	
	\begin{frame}{Q97: GDPR Rights - Tutorial}
		\fontsize{6.5}{7.5}\selectfont
		\textbf{QUESTION: What rights does GDPR give to data subjects?}
		
		\vspace{0.1cm}
		\textbf{OPTIONS:}
		\begin{enumerate}
			\item No rights
			\item The right to receive information, access data, request erasure, and object to processing
			\item Only the right to access data
			\item Only the right to delete data
		\end{enumerate}
		
		\vspace{0.1cm}
		\textcolor{myblue}{\textbf{ANSWER: B}}
		
		\vspace{0.1cm}
		\textbf{DETAILED EXPLANATION:}
		\begin{itemize}
			\item \textbf{Why B is correct:} Under GDPR, data subjects have rights to receive concise information, access personal data, request erasure, and object to processing for marketing or other purposes.
			\item \textbf{Why A is wrong:} GDPR gives significant rights
			\item \textbf{Why C is wrong:} It includes more than access
			\item \textbf{Why D is wrong:} It includes more than deletion
		\end{itemize}
		
		\textcolor{myred}{\textbf{EXAM TRAP:}} GDPR rights include \textcolor{myblue}{information, access, erasure, and objection}!
	\end{frame}
	
	\begin{frame}{Q98: GDPR Extraterritoriality - Tutorial}
		\fontsize{6.5}{7.5}\selectfont
		\textbf{QUESTION: Why has the GDPR been effective globally?}
		
		\vspace{0.1cm}
		\textbf{OPTIONS:}
		\begin{enumerate}
			\item Because it only applies in Europe
			\item Because it has extraterritorial jurisdiction, applying to companies outside Europe that offer services to European residents
			\item Because all countries adopted it
			\item Because it's optional
		\end{enumerate}
		
		\vspace{0.1cm}
		\textcolor{myblue}{\textbf{ANSWER: B}}
		
		\vspace{0.1cm}
		\textbf{DETAILED EXPLANATION:}
		\begin{itemize}
			\item \textbf{Why B is correct:} The GDPR applies not only if the data controller is inside Europe, but also if it's outside the European Economic Area if it offers services to data subjects within the EEA.
			\item \textbf{Why A is wrong:} It has extraterritorial effect
			\item \textbf{Why C is wrong:} Not all countries adopted it
			\item \textbf{Why D is wrong:} It's mandatory for those it applies to
		\end{itemize}
		
		\textcolor{myred}{\textbf{EXAM TRAP:}} GDPR has \textcolor{myblue}{extraterritorial jurisdiction} - applies globally to companies serving Europeans!
	\end{frame}
	
	\begin{frame}{Q99: EU AI Act - Tutorial}
		\fontsize{6.5}{7.5}\selectfont
		\textbf{QUESTION: What is the EU AI Act?}
		
		\vspace{0.1cm}
		\textbf{OPTIONS:}
		\begin{enumerate}
			\item A privacy regulation
			\item A law establishing a common regulatory and legal framework for AI
			\item A cybersecurity framework
			\item A voluntary guideline
		\end{enumerate}
		
		\vspace{0.1cm}
		\textcolor{myblue}{\textbf{ANSWER: B}}
		
		\vspace{0.1cm}
		\textbf{DETAILED EXPLANATION:}
		\begin{itemize}
			\item \textbf{Why B is correct:} The EU AI Act is a law establishing a common regulatory and legal framework for AI, approved by the EU Parliament in March 2024 and by the Council in May 2024.
			\item \textbf{Why A is wrong:} It's about AI, not just privacy
			\item \textbf{Why C is wrong:} It's about AI regulation, not just cybersecurity
			\item \textbf{Why D is wrong:} It's legally binding, not voluntary
		\end{itemize}
		
		\textcolor{myred}{\textbf{EXAM TRAP:}} EU AI Act = \textcolor{myblue}{common regulatory framework for AI}!
	\end{frame}
	
	\begin{frame}{Q100: Responsible AI Summary - Tutorial}
		\fontsize{6.5}{7.5}\selectfont
		\textbf{QUESTION: What is the most comprehensive summary of responsible AI?}
		
		\vspace{0.1cm}
		\textbf{OPTIONS:}
		\begin{enumerate}
			\item Only technical solutions matter
			\item Responsible AI requires ethical frameworks, privacy protection, bias mitigation, explainability, and robust governance and regulation
			\item Only regulation matters
			\item Ethics is optional
		\end{enumerate}
		
		\vspace{0.1cm}
		\textcolor{myblue}{\textbf{ANSWER: B}}
		
		\vspace{0.1cm}
		\textbf{DETAILED EXPLANATION:}
		\begin{itemize}
			\item \textbf{Why B is correct:} Responsible AI requires a comprehensive approach including ethical frameworks (consequentialism, deontology, virtue ethics), privacy protection, bias mitigation, explainability, and robust governance and regulation.
			\item \textbf{Why A is wrong:} Technical solutions alone are insufficient
			\item \textbf{Why C is wrong:} Ethics and other factors also matter
			\item \textbf{Why D is wrong:} Ethics is essential, not optional
		\end{itemize}
		
		\textcolor{myred}{\textbf{EXAM SUCCESS:}} Understand \textcolor{myblue}{ethical frameworks, principles, governance, and regulation} for responsible AI!
	\end{frame}
	
\end{document}